\documentclass[12pt,a4paper,UTF8]{ctexart}
\usepackage[a4paper,top=22mm,bottom=22mm,left=22mm,right=22mm,headheight=15pt]{geometry}
\newcommand{\Needspace}[1]{\par\begingroup\dimen0=#1\relax\dimen2=\pagegoal\advance\dimen2 by-\pagetotal\ifdim\dimen0>\dimen2\newpage\fi\endgroup}
\newcommand{\needspace}[1]{\par\begingroup\skip0=\lastskip\vskip-\skip0\dimen0=#1\relax\advance\dimen0 by\skip0\vskip\dimen0\penalty-100\vskip-\dimen0\vskip\skip0\penalty9999\endgroup}
\usepackage{amsmath,amssymb,mathtools,bm}
\usepackage{graphicx,xcolor,booktabs,longtable,array,enumitem,fancyhdr,caption,float,microtype,setspace}
\usepackage[hidelinks,unicode]{hyperref}
\usepackage[most]{tcolorbox}
\definecolor{ContestBlue}{HTML}{17365D}
\hypersetup{pdftitle={第070套},pdfauthor={}}
\setstretch{1.24}
\setlength{\parindent}{2em}
\setlength{\parskip}{0.35em}
\setlength{\emergencystretch}{2em}
\setlist{itemsep=0.35em,topsep=0.4em}
\captionsetup{font=small,labelfont=bf,justification=centering}
\setcounter{secnumdepth}{0}
\pagestyle{fancy}
\fancyhf{}
\fancyhead[L]{\small 第070套}
\fancyfoot[C]{\thepage}
\renewcommand{\headrulewidth}{0.4pt}
\ctexset{
  section={format=\Large\bfseries\color{ContestBlue},beforeskip=1.2em,afterskip=0.7em},
  subsection={format=\large\bfseries\color{ContestBlue},beforeskip=1.0em,afterskip=0.5em},
  subsubsection={format=\normalsize\bfseries,beforeskip=0.8em,afterskip=0.35em}
}
\usepackage{calc}
\newcounter{none}
\let\originalsubsection\subsection
\renewcommand{\subsection}{\needspace{6\baselineskip}\originalsubsection}
\makeatletter
\newcommand{\tightlist}{\setlength{\itemsep}{0pt}\setlength{\parskip}{0pt}}
\newcommand{\maxwidth}{\ifdim\Gin@nat@width>\linewidth\linewidth\else\Gin@nat@width\fi}
\newcommand{\maxheight}{\ifdim\Gin@nat@height>0.72\textheight 0.72\textheight\else\Gin@nat@height\fi}
\newcommand{\pandocbounded}[1]{\begingroup\setkeys{Gin}{width=\maxwidth,height=\maxheight,keepaspectratio}#1\endgroup}
\makeatother
\begin{document}
\begin{titlepage}
  \centering
  \vspace*{58mm}
  {\Huge\bfseries\color{ContestBlue} 第070套\par}
  \vfill
  \vspace*{28mm}
\end{titlepage}
\section{第 1 题}\label{ux7b2c-1-ux9898}

水平地面上有一足够长、质量为 \(m\) 的厚平板。平板侧边装有长度为 \(2l\)
的光滑发射台，发射方向与竖直方向夹角为 \(\theta\)。质量为 \(m\)
的小物块由理想弹簧推动；弹簧劲度系数

\[
k=\frac{2mg}{l},
\]

原长为 \(2l\)，压缩至长度 \(l\)
后用细线固定。系统初始静止，剪断细线后小物块沿发射台离开，并最终落回平板上，发射距离记为
\(s\)。

（1）平板固定时，求最大发射距离 \(s_1\) 与相应角度 \(\theta_1\)。

（2）平板可在光滑地面上运动时，最优角为 \(\theta_2\)。判断
\(\theta_2-\theta_1\) 的正负。

\begin{figure}[H]
\centering
\includegraphics[width=0.8\linewidth,height=\textheight,keepaspectratio,alt={图1-1　弹簧发射台与厚平板}]{../figures/question-01.pdf}
\caption*{图1-1　弹簧发射台与厚平板}
\end{figure}

\subsection{一、题型定位}\label{ux4e00ux9898ux578bux5b9aux4f4d}

弹簧推动物块上升，离台速度随发射角改变。固定板先求同高度射程的最大值；可动板再用水平动量守恒求反冲，并计算物块相对平板的落点。

\subsection{二、知识点}\label{ux4e8cux77e5ux8bc6ux70b9}

\begin{enumerate}
\def\labelenumi{\arabic{enumi}.}
\tightlist
\item
  弹簧势能与重力势能转化。
\item
  同高度抛体射程。
\item
  一元函数极值。
\item
  无水平外力系统的动量守恒。
\item
  相对载体的落点与飞行时间。
\end{enumerate}

\subsection{三、物理图像}\label{ux4e09ux7269ux7406ux56feux50cf}

增大 \(\theta\) 会减小抬升高度
\(l\cos\theta\)，使离台速度增大，同时减小竖直分速度所占比例。平板可动时，一部分能量转化为平板反冲，但相对平板的水平分速度仍由相对发射速度决定，最优角因此向更大的
\(\theta\) 偏移。

\subsection{四、解题入口}\label{ux56dbux89e3ux9898ux5165ux53e3}

先注意 \(\theta\) 是发射方向与竖直方向的夹角。弹簧从长度 \(l\) 恢复到
\(2l\) 时，物块沿台面移动 \(l\)，升高
\(l\cos\theta\)。第二问应把物块速度写成``相对平板速度 + 平板速度''。

\subsection{五、解析}\label{ux4e94ux89e3ux6790}

\begin{figure}[H]
\centering
\includegraphics[width=0.9\linewidth,height=\textheight,keepaspectratio,alt={图1-2　固定板与可动板的能量、动量和相对射程}]{../figures/solution-01.pdf}
\caption*{图1-2　固定板与可动板的能量、动量和相对射程}
\end{figure}

弹簧初始压缩量

\[
\Delta l=2l-l=l.
\tag{1}
\]

初始弹性势能

\[
E_s=\frac12kl^2
=\frac12\frac{2mg}{l}l^2
=mgl.
\tag{2}
\]

\needspace{3\baselineskip}
\textbf{（一）平板固定}

设物块离开发射台时速度为 \(v\)。在弹簧恢复过程中，物块升高

\[
\Delta h=l\cos\theta.
\tag{3}
\]

机械能守恒：

\[
\frac12mv^2+mgl\cos\theta=mgl.
\tag{4}
\]

所以

\[
\boxed{
v^2=2gl(1-\cos\theta)
}.
\tag{5}
\]

发射方向与竖直夹角为 \(\theta\)，故水平、竖直分速度大小分别为
\(v\sin\theta\)、\(v\cos\theta\)。飞行起点和落点等高，飞行时间

\[
T=\frac{2v\cos\theta}{g}.
\tag{6}
\]

射程

\[
\begin{aligned}
s_1(\theta)
&=v\sin\theta\,T\\
&=\frac{2v^2\sin\theta\cos\theta}{g}\\
&=4l\sin\theta\cos\theta(1-\cos\theta).
\end{aligned}
\tag{7}
\]

令

\[
x=\cos\theta,
\qquad 0<x<1,
\tag{8}
\]

则

\[
s_1=4l\,x(1-x)\sqrt{1-x^2}.
\tag{9}
\]

令 \(f(x)=x(1-x)\sqrt{1-x^2}\)。在 \(0<x<1\) 内
\(f(x)>0\)，可以对数求导：

\[
\frac{f'(x)}{f(x)}
=\frac1x-\frac1{1-x}-\frac{x}{1-x^2}.
\]

令 \(f'(x)=0\)，并乘以正量 \(x(1-x)(1+x)\)：

\[
(1-x^2)-x(1+x)-x^2=0.
\]

逐项合并后，极值条件化为

\[
3x^2+x-1=0.
\tag{10}
\]

方程的两根为 \(x=(-1\pm\sqrt{13})/6\)。其中负根不在 \(0<x<1\)
内，故物理解为

\[
x=\frac{\sqrt{13}-1}{6}.
\tag{11}
\]

因此

\[
\boxed{
\cos\theta_1=\frac{\sqrt{13}-1}{6}
},
\qquad
\boxed{\theta_1\approx64.3^\circ}.
\tag{12}
\]

最大射程

\[
\boxed{
s_{1,\max}
=4l\,x(1-x)\sqrt{1-x^2}
\bigg|_{x=(\sqrt{13}-1)/6}
\approx0.885\,l
}.
\tag{13}
\]

由于
\(f(0)=f(1)=0\)，而开区间内只有上述一个驻点，所以它一定给出全局最大值，而不是局部最小值。

\needspace{3\baselineskip}
\textbf{（二）平板可动}

设平板向左速度为 \(u\)，物块相对平板离台速度为
\(v\)。相对速度的水平分量为
\(v\sin\theta\)，故物块在地面系中的水平速度为

\[
v_x=v\sin\theta-u,
\tag{14}
\]

竖直速度仍为

\[
v_z=v\cos\theta.
\tag{15}
\]

水平方向动量守恒：

\[
m(v\sin\theta-u)-mu=0.
\tag{16}
\]

所以

\[
\boxed{u=\frac12v\sin\theta}.
\tag{17}
\]

能量守恒：

\[
\frac12mu^2
+\frac12m\left[(v\sin\theta-u)^2+v^2\cos^2\theta\right]
+mgl\cos\theta
=mgl.
\tag{18}
\]

代入式（17）并化简：

\[
\frac14mv^2(1+\cos^2\theta)
=mgl(1-\cos\theta).
\tag{19}
\]

故

\[
\boxed{
v^2=\frac{4gl(1-\cos\theta)}{1+\cos^2\theta}
}.
\tag{20}
\]

飞行期间平板保持匀速，物块相对平板的水平速度恰为
\(v\sin\theta\)。相对射程

\[
\begin{aligned}
s_2(\theta)
&=v\sin\theta\cdot\frac{2v\cos\theta}{g}\\
&=8l\frac{\sin\theta\cos\theta(1-\cos\theta)}
{1+\cos^2\theta}.
\end{aligned}
\tag{21}
\]

把固定板的公共因子记为

\[
f(\theta)=\sin\theta\cos\theta(1-\cos\theta).
\tag{22}
\]

则

\[
s_1=4lf(\theta),
\qquad
s_2=\frac{8lf(\theta)}{1+\cos^2\theta}.
\tag{23}
\]

在 \(\theta=\theta_1\) 处，\(f'(\theta_1)=0\)。对 \(s_2\) 求导：

\[
\left.\frac{ds_2}{d\theta}\right|_{\theta_1}
=8lf(\theta_1)
\left.
\frac{2\sin\theta\cos\theta}
{(1+\cos^2\theta)^2}
\right|_{\theta_1}>0.
\tag{24}
\]

因此 \(s_2\) 在 \(\theta_1\) 处仍随 \(\theta\) 增大，极大点位于其右侧：

\[
\boxed{\theta_2-\theta_1>0}.
\tag{25}
\]

\needspace{3\baselineskip}
\textbf{速度、落点和发射角放到同一图中}

图1-1的发射台长度为 \(2l\)，初始物块在中点，沿台面前进 \(l\) 后从平板
上沿离开。\(\theta\) 从竖直线量起，因此较大的 \(\theta\)
对应较平的发射方向。
图1-2改用离台时平板上沿为高度零点，抛物线在同一高度落回平板；
不能把落点画到地面后仍套用式（6）的等高飞行时间。

飞行期间，平板在地面系中向左移动 \(uT\)。物块地面系的水平位移为
\((v\sin\theta-u)T\)，相对平板的位移为

\[
s=(v\sin\theta-u)T+uT=v\sin\theta\,T.
\]

这一步直接对应图中的三个水平箭头。以相同角度比较两种模型，
\(s_2/s_1=2/(1+\cos^2\theta)\)；其范围从1到2。可动板并非把射程统一乘
一个常数，所以最优角才会改变。

令 \(x=\cos\theta\)，可动板驻点方程 \(p(x)=x^4+x^3+4x^2+x-1=0\) 在
\((0,1)\) 上严格递增。 在固定板根 \(3x_1^2+x_1-1=0\) 处，

\[
p(x_1)=x_1^2(x_1^2+x_1+1)>0.
\]

又 \(p(0)<0\)，故 \(0<x_2<x_1\)，从而 \(\theta_2>\theta_1\)。
这条代数比较不需要求四次方程的显式根，并补足``导数在旧极值处为正''
到``新全局极值在其右侧''之间的唯一性依据。

\needspace{3\baselineskip}
\textbf{补充推导与适用条件}

\textbf{无量纲结构。} 两种射程都可以写成 \(l\)
乘一个只含角度的无量纲函数，所以最优角与 \(m\)、\(g\)、\(l\)
的具体数值无关。\(m\) 在弹簧能、重力势能和动能中同时消去；\(g\)
决定离台速度和飞行时间，却在射程组合中约掉。这是检查推导是否合理的第一条线索。

\textbf{第二问的独立数值复核。} 若进一步令
\(x=\cos\theta\)，则可动板射程除去正常数后为

\[
g(x)=\frac{x(1-x)\sqrt{1-x^2}}{1+x^2}.
\]

在 \(0<x<1\) 内对数求导，驻点方程为

\[
x^4+x^3+4x^2+x-1=0.
\]

唯一正根为

\[
x_2=0.393526\ldots,
\qquad
\theta_2=66.82^\circ\ldots.
\]

固定板结果为
\(x_1=0.434259\ldots\)、\(\theta_1=64.26^\circ\ldots\)，确有
\(x_2<x_1\)、\(\theta_2>\theta_1\)。正文的导数符号法已经足够作答，这里计算精确驻点只是用来独立验证方向判断。

\textbf{端点解释。} 当 \(\theta\to0\)
时，物块几乎竖直发射，水平分速度趋于零；当 \(\theta\to90^\circ\)
时，弹簧不再克服重力，但竖直分速度趋于零。两端射程都为零，最大值必出现在中间。这个物理判断应先于求导出现。

\textbf{极值唯一性的代数核验。} 固定板的驻点方程为 \(3x^2+x-1=0\)。在
\(0<x<1\) 上，左端函数的导数 \(6x+1\) 恒正，且其在两端分别为 \(-1\) 和
\(3\)，因而区间内恰有一个根。结合两端射程为零，这个驻点必是全局最大值，而非局部最小值。可动板的四次多项式

\[
p(x)=x^4+x^3+4x^2+x-1
\]

也有 \(p'(x)=4x^3+3x^2+8x+1>0\)，同样在 \((0,1)\)
中只有一个根。这两次单调性检查补足了``解出一个驻点就宣称最大''所缺的逻辑环节。

\subsection{六、关键洞察}\label{ux516dux5173ux952eux6d1eux5bdf}

可动平板不需要显式求
\(\theta_2\)；只要在固定板最优点检查新射程函数的导数符号，就能判断最优角移动方向。

\subsection{七、易错点}\label{ux4e03ux6613ux9519ux70b9}

\begin{enumerate}
\def\labelenumi{\arabic{enumi}.}
\tightlist
\item
  把 \(\theta\) 当成与水平夹角。
\item
  弹簧能量全部转化为动能，漏掉发射台抬升。
\item
  可动板中直接套固定板速度。
\item
  用物块地面系水平速度计算相对板落点。
\item
  第二问耗费大量代数求精确 \(\theta_2\)，而题目只问正负。
\item
  把图中的 \(2l\) 看成平板长度；它是光滑发射台的总长度，平板本身足够长。
\item
  忘记弹簧恢复到原长时物块才离台，误把弹簧刚剪断处作为抛射起点。
\item
  在可动板能量式中只写平板动能与物块相对动能；动能必须在同一惯性系中计算。
\item
  由 \(x=\cos\theta\) 求出极值根后忘记 \(0<x<1\) 的范围，保留了负根。
\item
  用 \(ds/dx\) 的符号直接判断 \(ds/d\theta\)，忽略
  \(dx/d\theta=-\sin\theta<0\)。
\end{enumerate}

\subsection{八、课堂讲解建议}\label{ux516bux8bfeux5802ux8bb2ux89e3ux5efaux8bae}

先让学生在图 1-1 上指出两个长度为 \(l\) 的区段，并用手势明确 \(\theta\)
是与竖直方向的夹角。第一轮只处理离台过程：写出弹簧能、升高量与离台速度，不提前讨论抛体。

第二轮把离台点作为新的初态，列出
\(v_x=v\sin\theta\)、\(v_y=v\cos\theta\)，由竖直运动先求时间，再乘水平速度求射程。求导前要求学生先说明两个端点射程都为零，建立极值一定在内部的直觉。

第三轮让平板松开。先画速度三角形：相对速度为
\(\boldsymbol v\)，平板地面系速度为 \(-\boldsymbol u\)，物块地面系速度为
\(\boldsymbol v-\boldsymbol u\)。随后分别列水平动量和总能量。最后只在
\(\theta_1\) 处判断新函数导数，不急于求
\(\theta_2\)，突出``比较最优点移动方向''的方法。

课堂收尾可把数值
\(\theta_1\approx64.3^\circ\)、\(\theta_2\approx66.8^\circ\)
放回图中比较：在相同角度下，可动板的相对射程比固定板大，放大因子为
\(2/(1+\cos^2\theta)\)，它随 \(\theta\)
增大。因此最优方向向水平一侧移动。

\subsection{九、变式训练}\label{ux4e5dux53d8ux5f0fux8badux7ec3}

将平板质量改为 \(Mm\)。水平动量守恒给出

\[
u=\frac{v\sin\theta}{M+1}.
\]

把物块与平板的地面系动能相加，可得

\[
v^2=\frac{2gl(M+1)(1-\cos\theta)}{M+\cos^2\theta}.
\]

因此相对平板射程为

\[
\boxed{
s_M(\theta)=
\frac{4l(M+1)\sin\theta\cos\theta(1-\cos\theta)}
{M+\cos^2\theta}
}.
\]

令 \(M=1\) 恢复正文第二问；令 \(M\to\infty\)
恢复固定板射程。可进一步数值研究最优角随 \(M\) 增大如何连续回到
\(\theta_1\)。另可把弹簧劲度写成一般值 \(k\)，判断只有无量纲组合
\(kl/(mg)\) 会影响最优角。

对后一种变式可进一步完成计算。仍设弹簧压缩量和恢复过程的位移都为
\(l\)，定义

\[
\beta=\frac{kl}{2mg}.
\]

固定板时的能量方程为

\[
\frac12mv^2+mgl\cos\theta
=\frac12kl^2,
\]

故

\[
v^2=2gl(\beta-\cos\theta).
\]

射程为

\[
\boxed{
s(\theta)=4l\sin\theta\cos\theta
(\beta-\cos\theta)
}.
\]

物块能到达离台点还要求 \(\beta\ge\cos\theta\)。原题 \(k=2mg/l\) 对应
\(\beta=1\)，可行角度覆盖整个 \(0\le\theta\le90^\circ\)；若
\(\beta<1\)，较靠近竖直的发射方向会因弹簧能不足而变得不可达。这个可行域条件必须在求极值前先处理，否则可能得到数学驻点却没有对应的实际运动。

\newpage

\section{第 2 题}\label{ux7b2c-2-ux9898}

空间充满相对介电常数为 \(\varepsilon_{r2}\) 的均匀介质，远处匀强电场为

\[
\boldsymbol E_\infty=(E_0,0,E_0).
\]

在介质中挖出半径为 \(r\)、轴线沿 \(z\)
轴的无限长圆柱形孔，并在孔内填入相对介电常数为 \(\varepsilon_{r1}\)
的介质。求电场分布。

\begin{figure}[H]
\centering
\includegraphics[width=0.487\linewidth,height=\textheight,keepaspectratio,alt={图2-1　介质圆柱嵌入体与斜向外电场}]{../figures/question-02.pdf}
\caption*{图2-1　介质圆柱嵌入体与斜向外电场}
\end{figure}

\subsection{一、题型定位}\label{ux4e00ux9898ux578bux5b9aux4f4d-1}

圆柱轴向电场平行于介质界面，可直接保留；横向电场使界面极化。用柱内、柱外的一阶角谐波势函数，再由切向场和法向电位移的连续条件确定系数。

\subsection{二、知识点}\label{ux4e8cux77e5ux8bc6ux70b9-1}

\begin{enumerate}
\def\labelenumi{\arabic{enumi}.}
\tightlist
\item
  静电势满足拉普拉斯方程。
\item
  圆柱坐标的一阶角谐波解。
\item
  介质界面切向 \(E\) 连续、法向 \(D\) 连续。
\item
  叠加原理。
\item
  沿无限圆柱轴的电场分量不受界面扰动。
\end{enumerate}

\subsection{三、物理图像}\label{ux4e09ux7269ux7406ux56feux50cf-1}

外场的 \(z\) 分量始终平行于圆柱界面，边界条件自动满足，因此处处仍为
\(E_0\boldsymbol e_z\)。只有 \(x\)
分量使界面极化，产生二维偶极型扰动场。

\subsection{四、解题入口}\label{ux56dbux89e3ux9898ux5165ux53e3-1}

把问题拆成 \(E_0\boldsymbol e_x\) 与 \(E_0\boldsymbol e_z\)
两部分。横向势函数只需保留 \(\cos\phi\) 角谐波：柱内与 \(\rho\)
成正比，柱外由匀强场项和 \(1/\rho\) 偶极项组成。

\subsection{五、解析}\label{ux4e94ux89e3ux6790-1}

\begin{figure}[H]
\centering
\includegraphics[width=0.64\linewidth,height=\textheight,keepaspectratio,alt={图2-2　柱内均匀场与柱外偶极扰动场}]{../figures/solution-02.pdf}
\caption*{图2-2　柱内均匀场与柱外偶极扰动场}
\end{figure}

记

\[
\varepsilon_1=\varepsilon_0\varepsilon_{r1},
\qquad
\varepsilon_2=\varepsilon_0\varepsilon_{r2}.
\tag{1}
\]

\needspace{3\baselineskip}
\textbf{（一）轴向分量}

取

\[
\Phi_z=-E_0z.
\tag{2}
\]

它在界面内外连续，且法向电位移分量为零，因此

\[
\boxed{\boldsymbol E_z=E_0\boldsymbol e_z}
\tag{3}
\]

在全空间成立。

\needspace{3\baselineskip}
\textbf{（二）横向分量的势函数}

以圆柱坐标 \((\rho,\phi,z)\) 表示，\(x=\rho\cos\phi\)。设柱内势

\[
\Phi_{\rm in}^{(x)}
=-C E_0\rho\cos\phi,
\qquad \rho<r,
\tag{4}
\]

柱外势

\[
\Phi_{\rm out}^{(x)}
=-E_0\rho\cos\phi
+A\frac{\cos\phi}{\rho},
\qquad \rho>r.
\tag{5}
\]

在 \(\rho=r\)，势连续：

\[
-CE_0r=-E_0r+\frac Ar.
\tag{6}
\]

切向场因此也连续。法向位移连续：

\[
\varepsilon_1E_{\rho,\rm in}
=\varepsilon_2E_{\rho,\rm out}.
\tag{7}
\]

由势函数

\[
E_{\rho,\rm in}=CE_0\cos\phi,
\tag{8}
\]

\[
E_{\rho,\rm out}
=\left(E_0+\frac{A}{\rho^2}\right)\cos\phi.
\tag{9}
\]

由式（6）先把 \(A\) 用 \(C\) 表示：

\[
A=E_0r^2(1-C).
\]

在界面处，式（9）给出

\[
E_{\rho,\rm out}(r)
=E_0\left(1+\frac{A}{E_0r^2}\right)\cos\phi
=E_0(2-C)\cos\phi.
\]

代入法向位移连续条件并消去公共因子 \(E_0\cos\phi\)：

\[
\varepsilon_1C=\varepsilon_2(2-C),
\]

即

\[
(\varepsilon_1+\varepsilon_2)C=2\varepsilon_2.
\]

因此

\[
\boxed{
C=\frac{2\varepsilon_2}{\varepsilon_1+\varepsilon_2}
=\frac{2\varepsilon_{r2}}
{\varepsilon_{r1}+\varepsilon_{r2}}
},
\tag{10}
\]

\[
\boxed{
A=\frac{\varepsilon_1-\varepsilon_2}
{\varepsilon_1+\varepsilon_2}E_0r^2
}.
\tag{11}
\]

切向场也可直接复核。由两侧势函数

\[
E_{\phi,\rm in}=-CE_0\sin\phi,
\qquad
E_{\phi,\rm out}
=-\left(E_0-\frac{A}{\rho^2}\right)\sin\phi.
\]

令 \(\rho=r\) 并代入 \(A=E_0r^2(1-C)\)，外侧恰变为
\(-CE_0\sin\phi\)，所以切向连续没有遗漏。

\needspace{3\baselineskip}
\textbf{（三）最终电场}

柱内横向场均匀：

\[
\boldsymbol E_{\rm in}^{(x)}
=\frac{2\varepsilon_{r2}}
{\varepsilon_{r1}+\varepsilon_{r2}}
E_0\boldsymbol e_x.
\tag{12}
\]

加上轴向分量：

\[
\boxed{
\boldsymbol E_{\rm in}
=\frac{2\varepsilon_{r2}}
{\varepsilon_{r1}+\varepsilon_{r2}}
E_0\boldsymbol e_x
+E_0\boldsymbol e_z
}.
\tag{13}
\]

柱外定义

\[
\eta=
\frac{\varepsilon_{r1}-\varepsilon_{r2}}
{\varepsilon_{r1}+\varepsilon_{r2}}.
\tag{14}
\]

则

\[
\boxed{
E_{\rho,\rm out}
=E_0\left(1+\eta\frac{r^2}{\rho^2}\right)\cos\phi
},
\tag{15}
\]

\[
\boxed{
E_{\phi,\rm out}
=-E_0\left(1-\eta\frac{r^2}{\rho^2}\right)\sin\phi
},
\tag{16}
\]

\[
\boxed{E_{z,\rm out}=E_0}.
\tag{17}
\]

当 \(\varepsilon_{r1}=\varepsilon_{r2}\)
时，\(\eta=0\)，内外场均恢复为原匀强场。若
\(\varepsilon_{r1}\to\infty\)，柱内横向场趋于零，柱外
\(\eta\to1\)，又恢复二维导体圆柱在横向匀强场中的极限。这两个相互独立的极限共同核验了系数和符号。

\needspace{3\baselineskip}
\textbf{边界条件的逐分量检查}

图2-2在同一个界面点画出径向法线 \(\hat{\boldsymbol\rho}\) 与切向
\(\hat{\boldsymbol\phi}\)。应连续的是 \(E_\phi,E_z\) 和 \(D_\rho\)；
一般情况下 \(E_\rho\) 不连续。记
\(\eta=(\varepsilon_{r1}-\varepsilon_{r2})/
(\varepsilon_{r1}+\varepsilon_{r2})\)，则界面上

\[
E_{\rho,1}=CE_0\cos\phi,\quad E_{\rho,2}=(1+\eta)E_0\cos\phi,
\]

\[
E_{\phi,1}=E_{\phi,2}=-(1-\eta)E_0\sin\phi,\quad
C=1-\eta.
\]

\(\varepsilon_{r1}C=\varepsilon_{r2}(1+\eta)\) 保证法向电位移相等。
图中用 \(\varepsilon_{r1}=4,\varepsilon_{r2}=2\)
的横截面场箭头作具体示例， 柱内横场为远场的 \(2/3\)；轴向分量仍为
\(E_0\)，未画进二维截面。

在柱外把径向、切向场还原成笛卡尔分量，可写为

\[
E_x=E_0\left[1+\eta r^2\frac{x^2-y^2}{(x^2+y^2)^2}\right],\qquad
E_y=2E_0\eta r^2\frac{xy}{(x^2+y^2)^2}.
\]

这可用来绘制真实场箭头，并反查偶极项的符号。远处扰动按 \(\rho^{-2}\)
衰减。

将 \(\varepsilon_{r1}\to\infty\)
只作为横向介质响应的形式极限时，横场趋于
导体圆柱的二维结果；轴向场仍由原边值条件保留。不能把这整个含 \(E_z=E_0\)
的解称为静电平衡导体，因为导体内部总电场必须为零，
其轴向边界条件需要另行处理。

\needspace{3\baselineskip}
\textbf{补充推导与适用条件}

\textbf{为什么只保留一阶角谐波。} 远处横向匀强场的势为
\(-E_0\rho\cos\phi\)，其角向依赖只有
\(\cos\phi\)。圆柱界面是同心圆，边界条件不会把不同角谐波互相混合；因此柱内只需
\(\rho\cos\phi\)，柱外只需再加与之同阶的衰减解 \(\cos\phi/\rho\)。加入
\(\cos2\phi\) 等高阶项只会得到系数为零的齐次解。

\textbf{唯一性核验。}
已构造的势函数在两区域满足拉普拉斯方程，在无穷远处趋于给定匀强场，并在界面满足势连续与法向
\(D\)
连续。静电边值问题的唯一性定理保证它就是唯一解，而不只是``一个看起来合适的猜测''。

\textbf{极化方向。} 当 \(\varepsilon_{r1}>\varepsilon_{r2}\) 时
\(\eta>0\)，柱体沿外场方向更容易极化，柱内横向场被削弱；当
\(\varepsilon_{r1}<\varepsilon_{r2}\) 时
\(\eta<0\)，低介电常数夹杂内的横向场反而增强。轴向分量始终等于
\(E_0\)，所以柱内总场一般不再与远场方向平行。

\textbf{量纲检查。} 势函数中的 \(A\cos\phi/\rho\) 必须具有电势量纲，故
\(A\) 的量纲应为``场强乘长度平方''。式（11）恰为 \(E_0r^2\)
乘无量纲介电常数比；若写成 \(E_0r\) 或
\(E_0r^3\)，无需继续计算即可判错。

\textbf{界面极化电荷的反向核验。} 取单位法向
\(\hat{\boldsymbol n}=\hat{\boldsymbol\rho}\)
由柱内指向柱外。线性介质的极化强度为

\[
\boldsymbol P_i
=\varepsilon_0(\varepsilon_{ri}-1)\boldsymbol E_i.
\]

界面两侧极化的不连续给出净束缚面电荷

\[
\sigma_b=(\boldsymbol P_1-\boldsymbol P_2)\cdot\hat{\boldsymbol n}.
\]

由正文的界面场强

\[
E_{1n}=CE_0\cos\phi,
\qquad
E_{2n}=(1+\eta)E_0\cos\phi,
\]

以及 \(\varepsilon_{r1}C=\varepsilon_{r2}(1+\eta)\)，逐项化简得

\[
\begin{aligned}
\sigma_b
&=\varepsilon_0E_0\cos\phi
\left[(\varepsilon_{r1}-1)C
-(\varepsilon_{r2}-1)(1+\eta)\right]\\
&=\varepsilon_0E_0\cos\phi\,[(1+\eta)-C]\\
&=\boxed{2\varepsilon_0E_0\eta\cos\phi}.
\end{aligned}
\]

它在 \(+x\) 一侧与 \(-x\)
一侧符号相反，沿整个圆柱界面的总电荷为零，与``只极化、不凭空增加净电荷''的物理图像一致。这一结果也能反查法向边界条件的符号。

\subsection{六、关键洞察}\label{ux516dux5173ux952eux6d1eux5bdf-1}

三维问题因圆柱无限长而分解为``完全不受影响的轴向场''和``标准二维圆柱介质问题''。这种分解比直接在三维中猜极化电荷容易得多。

\subsection{七、易错点}\label{ux4e03ux6613ux9519ux70b9-1}

\begin{enumerate}
\def\labelenumi{\arabic{enumi}.}
\tightlist
\item
  把球形夹杂的系数 \(3\varepsilon_2/(\varepsilon_1+2\varepsilon_2)\)
  套到圆柱。
\item
  忘记外场还有 \(z\) 分量。
\item
  法向边界条件错误写成 \(E_n\) 连续。
\item
  外势偶极项写成 \(1/\rho^2\)；二维势应为 \(1/\rho\)。
\item
  未检验 \(\varepsilon_{r1}=\varepsilon_{r2}\) 极限。
\item
  在柱内势中保留 \(1/\rho\) 项，使轴线上出现无物理来源的奇点。
\item
  由势求 \(E_\phi\) 时漏掉圆柱坐标中的因子 \(-1/\rho\)。
\item
  把相对介电常数与绝对介电常数混用，却没有说明公共的 \(\varepsilon_0\)
  已约去。
\item
  求得柱内场后便停止，没有写出柱外扰动场，未完整回答``电场分布''。
\item
  忘记远场条件，给外势任意添加随 \(\rho\) 增长的齐次项。
\end{enumerate}

\subsection{八、课堂讲解建议}\label{ux516bux8bfeux5802ux8bb2ux89e3ux5efaux8bae-1}

先把外场箭头分解成平行柱轴和垂直柱轴两支。让学生只凭边界条件判断：平行分量没有法向位移，选择处处
\(E_z=E_0\) 已满足全部条件；需要求解的只剩横向问题。

横向部分先复习二维拉普拉斯方程的一阶角谐波。把柱内``在原点有限''和柱外``无穷远趋于匀强场''两个筛选条件写在黑板两侧，由学生判断为什么只剩

\[
\Phi_{\rm in}\propto\rho\cos\phi,
\qquad
\Phi_{\rm out}\propto-\rho\cos\phi+\frac{\cos\phi}{\rho}.
\]

再在图 2-2 的同一点 \(P\) 上画出
\(\hat{\boldsymbol\rho}\)、\(\hat{\boldsymbol\phi}\)，逐项计算
\(E_\rho,E_\phi\)。最后用``相同介质''\,``高介电常数柱''两个极限复核。这样的顺序能让边值问题从猜公式变成由正则性、远场和界面条件逐层筛选。

\subsection{九、变式训练}\label{ux4e5dux53d8ux5f0fux8badux7ec3-1}

对任意远场 \(\boldsymbol E_\infty\)，先分解

\[
\boldsymbol E_\parallel
=(\boldsymbol E_\infty\cdot\hat{\boldsymbol z})\hat{\boldsymbol z},
\qquad
\boldsymbol E_\perp
=\boldsymbol E_\infty-\boldsymbol E_\parallel.
\]

柱内统一结果为

\[
\boxed{
\boldsymbol E_{\rm in}
=\boldsymbol E_\parallel
+\frac{2\varepsilon_{r2}}
{\varepsilon_{r1}+\varepsilon_{r2}}
\boldsymbol E_\perp
}.
\]

空气孔只需令 \(\varepsilon_{r1}=1\)；外场完全沿轴时
\(\boldsymbol E_\perp=0\)，界面无扰动。还可把无限圆柱改为球形夹杂，比较球内系数
\(3\varepsilon_2/(\varepsilon_1+2\varepsilon_2)\)，由维数差异理解为何圆柱与球的去极化因子不同。

\newpage

\section{第 3 题}\label{ux7b2c-3-ux9898}

椭球壳内表面满足

\[
\frac{x^2}{a^2}+\frac{y^2}{b^2}+\frac{z^2}{c^2}=1.
\]

底部有半径为 \(r\)、质量为 \(m\)
的均匀实心小球，摩擦足够大，始终纯滚动。

（1）椭球壳固定。\(t=0\) 时球心位于最低点附近 \((0,y_0,z_0)\)，速度为
\((v_0,0,0)\)，并满足

\[
y_0\ll a,b,c,r,
\qquad
\frac{v_0^2}{g}\ll a,b,c,r.
\]

在小振幅近似下求 \(x(t)\)、\(y(t)\)。

（2）椭球壳以角速度 \(\omega_0\) 绕 \(z\)
轴匀速转动，求最低点附近小幅运动的简正频率。

\begin{figure}[H]
\centering
\includegraphics[width=0.636\linewidth,height=\textheight,keepaspectratio,alt={图3-1　椭球壳内滚动小球与三维坐标}]{../figures/question-03.pdf}
\caption*{图3-1　椭球壳内滚动小球与三维坐标}
\end{figure}

\subsection{一、题型定位}\label{ux4e00ux9898ux578bux5b9aux4f4d-2}

先沿法线把接触点平移到球心，求最低点附近球心轨迹的两个曲率。滚动惯量给出固定壳的频率；旋转壳还需规定轴向自转状态，才能确定速度耦合系数。

\subsection{二、知识点}\label{ux4e8cux77e5ux8bc6ux70b9-2}

\begin{enumerate}
\def\labelenumi{\arabic{enumi}.}
\tightlist
\item
  椭球最低点附近的二次曲面展开。
\item
  球心轨迹相对接触曲面的法向偏移。
\item
  实心球转动惯量 \(I=2mr^2/5\)。
\item
  纯滚动使有效惯性增为 \(m+I/r^2=7m/5\)。
\item
  旋转系中的离心、科里奥利力。
\item
  二维陀螺耦合特征方程。
\end{enumerate}

\subsection{三、物理图像}\label{ux4e09ux7269ux7406ux56feux50cf-2}

最低点附近，椭球内壁在 \(x\)、\(y\)
两个方向具有不同曲率，故固定壳时存在两个互相独立的简谐频率。壳旋转后，科里奥利力把
\(x\)、\(y\) 两方向耦合，正常模态变成椭圆偏振运动。

\subsection{四、解题入口}\label{ux56dbux89e3ux9898ux5165ux53e3-2}

先求球心有效轨迹的曲率，而不是直接把球心放在原椭球面上。固定壳时由纯滚动的有效质量得到两个频率；旋转壳时在随壳转动系中加入离心修正与速度耦合。

\subsection{五、解析}\label{ux4e94ux89e3ux6790-2}

\begin{figure}[H]
\centering
\includegraphics[width=0.9\linewidth,height=\textheight,keepaspectratio,alt={图3-2　局部曲率、独立振动与旋转系耦合}]{../figures/solution-03.pdf}
\caption*{图3-2　局部曲率、独立振动与旋转系耦合}
\end{figure}

\needspace{3\baselineskip}
\textbf{（一）固定椭球壳}

椭球最低点在 \(z=-c\)。对内壁接触点
\((x',y',z')\)，最低点附近单位法向近似为

\[
\boldsymbol n
\approx
\frac{c}{a^2}x'\boldsymbol e_x
+\frac{c}{b^2}y'\boldsymbol e_y
-\boldsymbol e_z.
\tag{1}
\]

式（1）的 \(\boldsymbol n\) 取壳面向外法线，球心为接触点减去
\(r\boldsymbol n\)，即沿内法线偏移 \(r\)。因此球心横坐标

\[
x=x'-r\frac{c}{a^2}x'
=\frac{a^2-cr}{a^2}x',
\tag{2}
\]

\[
y=\frac{b^2-cr}{b^2}y'.
\tag{3}
\]

由椭球面二阶展开，球心高度相对最低点的增量为

\[
\Delta z
\approx
\frac{c}{2(a^2-cr)}x^2
+\frac{c}{2(b^2-cr)}y^2.
\tag{4}
\]

重力势能

\[
U
=\frac12k_xx^2+\frac12k_yy^2,
\tag{5}
\]

其中

\[
k_x=\frac{mgc}{a^2-cr},
\qquad
k_y=\frac{mgc}{b^2-cr}.
\tag{6}
\]

纯滚动时

\[
T=\frac12m(\dot x^2+\dot y^2)
+\frac12\frac{2}{5}mr^2
\frac{\dot x^2+\dot y^2}{r^2}.
\tag{7}
\]

所以

\[
T=\frac12\frac75m(\dot x^2+\dot y^2).
\tag{8}
\]

运动方程

\[
\ddot x+\omega_x^2x=0,
\qquad
\ddot y+\omega_y^2y=0,
\tag{9}
\]

其中

\[
\boxed{
\omega_x^2=\frac{5gc}{7(a^2-cr)}
},
\qquad
\boxed{
\omega_y^2=\frac{5gc}{7(b^2-cr)}
}.
\tag{10}
\]

由初始条件

\[
x(0)=0,\quad \dot x(0)=v_0,
\qquad
y(0)=y_0,\quad \dot y(0)=0,
\tag{11}
\]

得到

\[
\boxed{
x(t)=\frac{v_0}{\omega_x}\sin\omega_xt
},
\tag{12}
\]

\[
\boxed{
y(t)=y_0\cos\omega_yt
}.
\tag{13}
\]

\needspace{3\baselineskip}
\textbf{（二）椭球壳绕 \(z\) 轴旋转：参考解采用的自转状态}

此处 \(x,y\)
均在随壳旋转的坐标系中。参考解的滚动约束隐含球的绝对轴向自转分量
\(\Omega_z=\omega_0\)，即球相对壳无轴向自转。原题的``纯滚动''只约束接触点速度，本身未规定绕接触法线的自转；若允许任意轴向自转，频率还需这个参数。以下先按参考解的状态求解，再给出一般式。

在随壳转动系中，离心势使两个方向的有效劲度变为

\[
k_x'=k_x-m\omega_0^2,
\qquad
k_y'=k_y-m\omega_0^2.
\tag{14}
\]

纯滚动的平动---转动耦合使科里奥利项的有效系数为

\[
q=\frac{12}{7}\omega_0.
\tag{15}
\]

定义

\[
A=\frac{5k_x'}{7m}
=\frac{5gc}{7(a^2-cr)}-\frac57\omega_0^2,
\tag{16}
\]

\[
B=\frac{5k_y'}{7m}
=\frac{5gc}{7(b^2-cr)}-\frac57\omega_0^2.
\tag{17}
\]

下面补出式（15）的来龙去脉。记实心球转动惯量

\[
I=\kappa mr^2,
\qquad
\kappa=\frac25.
\]

在随壳转动系中，水平平动方程为

\[
\begin{aligned}
m\ddot x
&=-k_xx+f_x+2m\omega_0\dot y+m\omega_0^2x,\\
m\ddot y
&=-k_yy+f_y-2m\omega_0\dot x+m\omega_0^2y.
\end{aligned}
\]

这里 \(f_x,f_y\) 是保证纯滚动的静摩擦分量；两项 \(\pm2m\omega_0\dot q\)
来自科里奥利力，两项 \(m\omega_0^2q\) 来自离心力。无滑动约束给出

\[
\omega_xr=-\dot y,
\qquad
\omega_yr=\dot x.
\]

球的自转角速度分量是在旋转基底中表示的，所以角动量方程必须保留基底转动项：

\[
I(\dot\omega_x-\omega_0\omega_y)=rf_y,
\qquad
I(\dot\omega_y+\omega_0\omega_x)=-rf_x.
\]

对无滑动约束求导，再代入上式：

\[
f_x=-\kappa m\ddot x+\kappa m\omega_0\dot y,
\qquad
f_y=-\kappa m\ddot y-\kappa m\omega_0\dot x.
\]

把摩擦分量代回平动方程，\(x\) 方向得到

\[
(1+\kappa)\ddot x
+\left(\frac{k_x}{m}-\omega_0^2\right)x
-(2+\kappa)\omega_0\dot y=0.
\]

除以 \(1+\kappa=7/5\)，速度耦合系数为

\[
\frac{2+\kappa}{1+\kappa}\omega_0
=\frac{12}{7}\omega_0=q.
\]

\(y\) 方程的推导相同，只是速度耦合项符号相反。于是线性化方程为

\[
\boxed{
\begin{cases}
\ddot x+A x-q\dot y=0,\\
\ddot y+B y+q\dot x=0.
\end{cases}
}
\tag{18}
\]

取

\[
x=Xe^{i\omega t},
\qquad
y=Ye^{i\omega t},
\tag{19}
\]

得到特征方程

\[
(A-\omega^2)(B-\omega^2)-q^2\omega^2=0.
\tag{20}
\]

展开为

\[
\omega^4-(A+B+q^2)\omega^2+AB=0.
\tag{21}
\]

因此两个简正频率为

\[
\boxed{
\omega_\pm^2
=\frac{
A+B+q^2
\pm\sqrt{(A+B+q^2)^2-4AB}
}{2}
}.
\tag{22}
\]

\(A>0\)、\(B>0\)
是稳定小振动的直接充分条件。更一般地，应检查式（22）的两个 \(\omega^2\)
是否均为正实数；陀螺耦合较强时，不能只凭某一个静态劲度的符号武断地下结论。

\needspace{3\baselineskip}
\textbf{球心平行曲面的二阶展开}

定义壳面最低点的两主曲率 \(\kappa_x=c/a^2\)、\(\kappa_y=c/b^2\)。
接触点高度为 \(z'=-c+\tfrac12\kappa_xx'^2+\tfrac12\kappa_yy'^2\)。
外法线的竖直分量为

\[
n_z=-1+\frac12(\kappa_x^2x'^2+\kappa_y^2y'^2).
\]

球心 \(C=P-r\boldsymbol n\)，所以

\[
x=(1-r\kappa_x)x',\quad y=(1-r\kappa_y)y',
\]

\[
z_C=-c+r+\frac12\kappa_x(1-r\kappa_x)x'^2
+\frac12\kappa_y(1-r\kappa_y)y'^2.
\]

消去接触坐标，便得到式（4）。要求 \(r<a^2/c\) 且
\(r<b^2/c\)，保证球能在底部以单点接触安放，
球心轨迹曲率有限且正。图3-2明确把接触点 \(P\) 与球心 \(C\) 分开， 线段
\(PC=r\) 沿法线，两个局部曲率半径分别减去 \(r\)。

当 \(q=0\) 且 \(A=B>0\)，两个频率相等，判别式为零，但仍是稳定的两个
独立振子；因此充分条件证明中的判别式只能写``非负''，不能要求严格正。

\needspace{3\baselineskip}
\textbf{任意轴向自转下的频率}

令 \(\Omega_s\) 为球在最低点的绝对轴向自转角速度，它在线性阶为常数， 记
\(w=\Omega_s-\omega_0\)、\(\beta=I/(mr^2)\)，并定义球心轨迹曲率

\[
K_x=\frac{c}{a^2-cr},\qquad K_y=\frac{c}{b^2-cr}.
\]

接触点速度相等给出

\[
\Omega_x=-\frac{\dot y}{r}-wK_xx,\qquad
\Omega_y=\frac{\dot x}{r}-wK_yy.
\]

把这两式代入旋转基底中的角动量方程，再消去摩擦，有

\[
\ddot x+A_sx-q_x\dot y=0,\qquad
\ddot y+B_sy+q_y\dot x=0,
\]

其中

\[
\begin{aligned}
A_s&=\frac{gK_x-\omega_0^2-\beta\omega_0rwK_x}{1+\beta},\\
B_s&=\frac{gK_y-\omega_0^2-\beta\omega_0rwK_y}{1+\beta},\\
q_x&=\frac{(2+\beta)\omega_0+\beta rwK_y}{1+\beta},\\
q_y&=\frac{(2+\beta)\omega_0+\beta rwK_x}{1+\beta}.
\end{aligned}
\]

频率方程为

\[
\omega^4-(A_s+B_s+q_xq_y)\omega^2+A_sB_s=0.
\]

当 \(w=0\)、\(\beta=2/5\)，立即回到正文的 \(A,B,q=12\omega_0/7\)。
若取实验室系中初始无轴向自转 \(\Omega_s=0\)，则 \(w=-\omega_0\)，
不能再使用同一组频率。图中的科里奥利箭头只按参考解的状态示意，
方向还取决于实际速度；它不是沿轨迹固定指向某一方向的回复力。

\needspace{3\baselineskip}
\textbf{补充推导与适用条件}

\textbf{两次修正不能混为一谈。} \(a^2-cr\)、\(b^2-cr\)
来自``球心轨迹不是接触点轨迹''的几何偏移；\(5/7\)
来自平动动能与自转动能之和。前者改变回复劲度，后者改变惯性。若先把球当质点求完频率再简单乘
\(\sqrt{5/7}\)，数值上虽可能碰巧一致，却会掩盖有限半径对曲率的另一重修正。

\textbf{固定壳的能量核验。} 对 \(x\) 方向单独写

\[
E_x=\frac12\frac75m\dot x^2
+\frac12\frac{mgc}{a^2-cr}x^2.
\]

求导并令 \(dE_x/dt=0\)，立即得到
\(\ddot x+[5gc/7(a^2-cr)]x=0\)。这与牛顿---转动方程消去摩擦所得结果一致，也说明静摩擦在纯滚动中不做功。

\textbf{轴对称极限。} 若 \(a=b\)，则
\(A=B\equiv\omega_r^2\)。两个圆偏振频率可从式（22）进一步写成

\[
\omega_\pm
=\sqrt{\omega_r^2+\frac{q^2}{4}}
\pm\frac q2.
\]

一个模式相对壳顺转，另一个逆转；\(\omega_0\to0\)
时二者同时回到固定壳频率。这个极限同时核对了特征行列式中 \(q^2\omega^2\)
的符号。

\textbf{适用范围。} 所有式子只保留 \(x,y\)
的一阶动力学和势能的二阶项，并假设接触处不失去纯滚动。若振幅大到曲率显著变化，或所需静摩擦超过上限，线性简正频率不再适用。

\textbf{有效拉格朗日量的独立复核。}
将拉格朗日量除去共同有效质量后，旋转壳中的两个线性方程可由

\[
L=\frac12(\dot x^2+\dot y^2)
+\frac q2(x\dot y-y\dot x)
-\frac12(Ax^2+By^2)
\]

给出。对 \(x\) 使用欧拉---拉格朗日方程：

\[
\frac{d}{dt}\frac{\partial L}{\partial\dot x}
-\frac{\partial L}{\partial x}
=(\ddot x-\tfrac q2\dot y)
-(\tfrac q2\dot y-Ax)=0,
\]

即 \(\ddot x-q\dot y+Ax=0\)；对 \(y\) 同理得
\(\ddot y+q\dot x+By=0\)。其中含速度的项与带电粒子在磁场中的矢势项形式相同，它只改变运动方向，不出现在守恒能量

\[
E=\frac12(\dot x^2+\dot y^2)
+\frac12(Ax^2+By^2)
\]

中。这个结构直观地说明科里奥利耦合本身不做功。

\textbf{稳定性判据为何成立。} 令 \(z=\omega^2\)，特征方程成为

\[
z^2-(A+B+q^2)z+AB=0.
\]

若 \(A>0,B>0\)，两根之和为正、两根之积为正。其判别式还可改写为

\[
\Delta=(A-B)^2+2q^2(A+B)+q^4\ge0.
\]

因此两个 \(z\)
都是正实数，两个模态频率都为实数。这一步把正文中``\(A,B>0\)
为充分条件''展开成了可逐项检查的代数证明。

\textbf{模态为何是椭圆运动。} 把某一本征频率 \(\omega\)
代回第一行振幅方程，得

\[
(A-\omega^2)X-iq\omega Y=0,
\]

即

\[
\frac YX
=-i\frac{A-\omega^2}{q\omega}.
\]

右端是纯虚数，说明 \(x\) 与 \(y\) 相差
\(\pi/2\)；其绝对值给出椭圆两个半轴振幅之比。例如取 \(X\)
为实数，则实运动可写成

\[
x(t)=X\cos\omega t,
\qquad
y(t)=\frac{A-\omega^2}{q\omega}X\sin\omega t.
\]

系数的正负还决定该模态在 \(x\)---\(y\) 平面中顺时针还是逆时针旋转。当
\(A=B\) 时，利用圆模态频率关系可得振幅比绝对值为 \(1\)，椭圆退化为圆；当
\(q\to0\) 时，两个模态则分别退化为沿 \(x\) 或 \(y\) 的直线振动。

\needspace{3\baselineskip}
\textbf{一般自转状态下的摩擦力消去}

前面给出的 \(\Omega_{{\rm rel},x},\Omega_{{\rm rel},y}\)
在随壳旋转的基底中表示。实验室角加速度必须包括基底转动：

\[
\alpha_x=\dot\Omega_x-\omega_0\Omega_y,\qquad
\alpha_y=\dot\Omega_y+\omega_0\Omega_x.
\]

接触摩擦关于球心的水平力矩为
\((rf_y,-rf_x)\)，法向支持力通过球心，不提供自转力矩。用
\(I=\beta mr^2\) 并保留一阶量，得到

\[
\begin{aligned}
f_x&=-\beta m\ddot x+\beta m(\omega_0+rwK_y)\dot y
 +\beta m\omega_0rwK_xx,\\
f_y&=-\beta m\ddot y-\beta m(\omega_0+rwK_x)\dot x
 +\beta m\omega_0rwK_yy.
\end{aligned}
\]

将它们代入随壳坐标中的水平平动式

\[
\ddot x-2\omega_0\dot y-\omega_0^2x=-gK_xx+f_x/m,
\]

以及对应的 \(y\) 式，便恢复前面的
\(A_s,B_s,q_x,q_y\)。这个步骤显示，轴向自转既通过倾斜法线影响滚动约束，也通过旋转基底影响角加速度；漏去任何一项都不能靠修改最后一个系数补救。

取 \(w=0\)，摩擦式恢复参考分支；取 \(\omega_0=w=0\)，则
\(f_x=-\beta m\ddot x\)，有效平动惯量为
\((1+\beta)m\)。这两个检查分别检验旋转耦合与固定壳的滚动惯量。

\needspace{3\baselineskip}
\textbf{负静态劲度时的陀螺稳定与重根}

参考分支的特征式令 \(z=\omega^2\) 后为 \(z^2-\Sigma z+\Pi=0\)，其中
\(\Sigma=A+B+q^2,\Pi=AB\)。两根严格正且不同的充要代数条件是

\[
\Sigma>0,\qquad \Pi>0,\qquad
\Delta=\Sigma^2-4\Pi>0.
\]

若 \(A,B\) 异号，则
\(\Pi<0\)，一个平方频率必为负，速度耦合无法将它稳定。若 \(A,B\)
同为负，\(\Pi>0\)，足够强的陀螺耦合仍可能使 \(\Sigma,\Delta\)
满足上述条件。此时静态势能并非极小，但线性轨道可以有界。

例如取圆对称曲率并以 \(gK=1\)
为时间单位，\(\omega_0=3\)、\(w=0\)、\(\beta=2/5\)。有

\[
A=B=-40/7,\qquad q=36/7,
\]

\[
z_\pm=\frac{736\pm\sqrt{228096}}{98}>0.
\]

这例说明 \(A,B>0\)
是方便的充分条件，而不是所有旋转速度下的必要条件。它仍须满足小振幅、单点接触与静摩擦足以维持滚动；不能把线性谱的结论推广到有限振幅稳定。

判别式等于零时，还要检查本征矢量是否足够。对
\(A=B>0,q=0\)，系统是两个独立同频振子，重根没有问题；对
\(A=B=-a<0,q=2\sqrt a\)，圆偏振方程出现重复的时间根，一般解含
\(te^{\pm i\sqrt a t}\)，会缓慢无界增长。只检查 \(\omega^2\)
为正实数会漏掉这种Jordan退化。故严格稳定区与其临界边界应分开说明。

\subsection{六、关键洞察}\label{ux516dux5173ux952eux6d1eux5bdf-2}

半径 \(r\) 通过 \(a^2-cr\)、\(b^2-cr\)
修正球心轨迹曲率；纯滚动再给出统一的 \(5/7\)
系数。旋转后不仅有离心软化，还有不可忽略的科里奥利耦合。

\subsection{七、易错点}\label{ux4e03ux6613ux9519ux70b9-2}

\begin{enumerate}
\def\labelenumi{\arabic{enumi}.}
\tightlist
\item
  直接把球心坐标代入原椭球方程。
\item
  把小球当质点，漏掉 \(5/7\)。
\item
  旋转系只加离心力，漏掉速度耦合。
\item
  特征行列式中把 \(q^2\omega^2\) 符号写错。
\item
  不检查 \(A\)、\(B\) 的稳定性。
\item
  把球心轨迹半轴误写成
  \(a-r\)、\(b-r\)；平行曲面的局部曲率修正不是简单减半径。
\item
  认为静摩擦必等于
  \(\mu N\)。纯滚动阶段静摩擦由动力学反求，只需不超过上限。
\item
  在旋转基底中对自转分量直接求普通时间导数，漏掉
  \(\boldsymbol{\omega}_0\times\boldsymbol\omega\)。
\item
  用 \(x=X\cos\omega t\)、\(y=Y\cos\omega t\)
  强行取同相实振幅，无法表示椭圆偏振模态；应使用复振幅。
\item
  只检查 \(\omega_\pm^2\) 的代数值，没有核对 \(\omega_0\to0\) 与 \(a=b\)
  两个极限。
\end{enumerate}

\subsection{八、课堂讲解建议}\label{ux516bux8bfeux5802ux8bb2ux89e3ux5efaux8bae-2}

先在 \(x\)---\(z\) 截面画椭圆内壁、接触点和球心。让学生沿法线平移距离
\(r\)，只保留 \(x\) 的一阶关系与高度的二阶关系，由此得到
\(a^2-cr\)。再把同一过程复制到 \(y\)
方向，强调这是两个主曲率方向的独立局部展开。

随后只讨论固定壳。把平动动能和自转动能并排写出，让学生亲自得到有效质量
\(7m/5\)，再用能量法写两个独立振子。此处可用 \(r\to0\)
检查几何修正，但要提醒：即使
\(r\to0\)，若仍保留实心球滚动，转动惯性因子不会自动变成质点。

旋转部分分三层：先写离心软化，再写平动科里奥利力，最后说明滚动约束使球的自转分量也受旋转基底影响。将
\(f_x,f_y\) 消去后，才出现
\(12\omega_0/7\)。最后用``带磁场的二维各向异性振子''类比说明速度耦合不改变势能，却改变模态旋转方向。

求出特征方程后，让学生先检查 \(\omega_0=0\)，再令 \(a=b\)
写成圆偏振频率。两个极限都闭合后，再讨论一般稳定性。

\subsection{九、变式训练}\label{ux4e5dux53d8ux5f0fux8badux7ec3-2}

把小球换成转动惯量

\[
I=\kappa mr^2
\]

的一般滚动物体。固定壳频率变为

\[
\omega_x^2
=\frac{gc}{(1+\kappa)(a^2-cr)},
\qquad
\omega_y^2
=\frac{gc}{(1+\kappa)(b^2-cr)}.
\]

旋转时

\[
q=\frac{2+\kappa}{1+\kappa}\omega_0,
\]

而 \(A,B\) 的共同分母也由 \(7/5\) 改为 \(1+\kappa\)。代入薄环
\(\kappa=1\)、空心球 \(\kappa=2/3\) 与实心球
\(\kappa=2/5\)，比较转动惯量对频率与陀螺分裂的影响。另可令
\(a=b\)，验证顺、逆旋圆模态；若彻底改为无滚动质点，应同时令 \(\kappa=0\)
并删除无滑动约束。

\newpage

\section{第 4 题}\label{ux7b2c-4-ux9898}

如图所示，自制手推炮的炮膛截面积为 \(S\)，工作过程分为以下四个状态。

\begin{figure}[H]
\centering
\includegraphics[width=0.744\linewidth,height=\textheight,keepaspectratio,alt={图4-1　手推炮的四个工作状态与气柱几何量}]{../figures/question-04.pdf}
\caption*{图4-1　手推炮的四个工作状态与气柱几何量}
\end{figure}

\begin{enumerate}
\def\labelenumi{\arabic{enumi}.}
\tightlist
\item
  炮膛中已有一颗炮弹 \(A\)，再从炮尾塞入一颗与其相同的炮弹
  \(B\)。两弹质量均为
  \(m\)，体积均可忽略，且与炮膛之间保持良好气密性。随后以速度 \(u\)
  匀速推动手柄，压缩两弹之间的气体。
\item
  当炮弹 \(B\) 被推到距炮口 \(l_1+l_2\) 处时，炮弹 \(A\) 已运动到距炮口
  \(l_1\) 处，并获得速度 \(u\)。此时两弹间气体压强为
\end{enumerate}

\[
P_0=\frac{ma_0}{S}.
\]

此后 \(B\) 继续由手柄匀速推进，\(A\)
在气体压力作用下加速。作如下理想化：当炮弹相对炮膛的速度不小于 \(u\)
时，炮弹与炮膛之间的摩擦可忽略；气体经历绝热膨胀，且其压强始终远大于大气压；绝热指数为

\[
   \gamma=\frac{C_P}{C_V}=2.
   \]

\begin{enumerate}
\def\labelenumi{\arabic{enumi}.}
\setcounter{enumi}{2}
\tightlist
\item
  从状态 2 到状态 3，炮弹 \(B\) 被手柄匀速推进距离 \(h\)，炮弹 \(A\)
  被加速至速度 \(v\) 并射出。
\item
  发射完成后，炮弹 \(B\) 继续向炮口推进并成为下一次发射中的炮弹 \(A\)。
\end{enumerate}

（1）忽略状态 2 到状态 3 中炮弹 \(B\) 的推进距离 \(h\)，求炮弹 \(A\)
的发射速度 \(v_1\)，结果用 \(u,l_1,l_2,a_0\) 表示。

（2）计入推进距离 \(h\)，求炮弹 \(A\) 的发射速度 \(v_2\)，结果用
\(h,u,l_1,l_2,a_0\) 表示。

（3）根据炮弹的动力学过程确定 \(h\) 与 \(u,l_1,l_2,a_0\) 的关系。若
\(h=l\)、\(l_1=2l\)、\(l_2=3l\)，求 \(u\)，结果用 \(l,a_0\) 表示。

\subsection{一、题型定位}\label{ux4e00ux9898ux578bux5b9aux4f4d-3}

两弹间气体绝热膨胀推动前弹，后弹同时由手柄匀速推进。先区分前弹的炮膛位移与两弹的相对位移，再由压强做功求速度、由时间积分闭合手柄推进距离。

\subsection{二、知识点}\label{ux4e8cux77e5ux8bc6ux70b9-3}

\begin{enumerate}
\def\labelenumi{\arabic{enumi}.}
\tightlist
\item
  绝热关系 \(PV^\gamma=\text{常量}\)。
\item
  压强功与炮弹动能增量。
\item
  匀速移动活塞参考系。
\item
  由速度---位移关系积分求时间。
\item
  自洽条件 \(h=uT\)。
\end{enumerate}

\subsection{三、物理图像}\label{ux4e09ux7269ux7406ux56feux50cf-3}

状态 2 时两弹速度同为 \(u\)，在 \(B\) 的参考系中 \(A\) 初速度为零。随着
\(A\) 向炮口运动、\(B\) 同向推进，两弹间距的净增量为
\(l_1-h\)，气体膨胀推动 \(A\)。

\subsection{四、解题入口}\label{ux56dbux89e3ux9898ux5165ux53e3-3}

先写气体压强随两弹间距的变化。第一问在炮膛系做功；第二问切换到随 \(B\)
匀速运动的参考系，初始相对动能为零，最后再加回 \(u\)。

\subsection{五、解析}\label{ux4e94ux89e3ux6790-3}

\begin{figure}[H]
\centering
\includegraphics[width=0.885\linewidth,height=\textheight,keepaspectratio,alt={图4-2　气体压强曲线、相对位移与发射时间积分}]{../figures/solution-04.pdf}
\caption*{图4-2　气体压强曲线、相对位移与发射时间积分}
\end{figure}

初始气柱长度为 \(l_2\)。绝热指数 \(\gamma=2\)，故当气柱长度变为
\(l_2+x\) 时

\[
P(x)(l_2+x)^2=P_0l_2^2.
\tag{1}
\]

所以

\[
\boxed{
P(x)=P_0\frac{l_2^2}{(l_2+x)^2}
}.
\tag{2}
\]

\needspace{3\baselineskip}
\textbf{（一）忽略 \(B\) 的推进}

\(A\) 从距炮口 \(l_1\) 处运动到炮口，气柱伸长量从 \(x=0\) 增至
\(x=l_1\)。压强功

\[
W=\int_0^{l_1}P(x)S\,dx.
\tag{3}
\]

代入 \(P_0S=ma_0\)：

\[
\begin{aligned}
W
&=ma_0l_2^2\int_0^{l_1}\frac{dx}{(l_2+x)^2}\\
&=ma_0l_2^2
\left(\frac1{l_2}-\frac1{l_1+l_2}\right)\\
&=ma_0\frac{l_1l_2}{l_1+l_2}.
\end{aligned}
\tag{4}
\]

动能定理：

\[
\frac12m(v_1^2-u^2)=W.
\tag{5}
\]

因此

\[
\boxed{
v_1=\sqrt{
u^2+\frac{2a_0l_1l_2}{l_1+l_2}
}
}.
\tag{6}
\]

\needspace{3\baselineskip}
\textbf{（二）计入 \(B\) 的推进}

在随 \(B\) 以速度 \(u\) 运动的惯性系中，状态 2 时 \(A\)
的相对速度为零。到 \(A\) 出口时，\(A\) 前进 \(l_1\)，\(B\) 前进
\(h\)，故气柱净伸长

\[
x_f=l_1-h.
\tag{7}
\]

相对系动能定理：

\[
\frac12m v_{\rm rel}^2
=ma_0l_2^2
\int_0^{l_1-h}\frac{dx}{(l_2+x)^2}.
\tag{8}
\]

积分：

\[
\frac12m v_{\rm rel}^2
=ma_0\frac{(l_1-h)l_2}{l_1-h+l_2}.
\tag{9}
\]

故

\[
v_{\rm rel}
=\sqrt{
\frac{2a_0(l_1-h)l_2}
{l_1-h+l_2}
}.
\tag{10}
\]

回到炮膛系：

\[
\boxed{
v_2
=u+\sqrt{
\frac{2a_0(l_1-h)l_2}
{l_1-h+l_2}
}
}.
\tag{11}
\]

\needspace{3\baselineskip}
\textbf{（三）由运动时间确定推进速度}

在 \(B\) 系中，当相对位移为 \(x\) 时，由式（8）的中间结果

\[
v_{\rm rel}(x)
=\sqrt{\frac{2a_0l_2x}{l_2+x}}.
\tag{12}
\]

因此

\[
T=\int_0^{l_1-h}\frac{dx}{v_{\rm rel}(x)}
=\frac1{\sqrt{2a_0l_2}}
\int_0^{l_1-h}\sqrt{\frac{l_2+x}{x}}\,dx.
\tag{13}
\]

代入

\[
h=l,\qquad l_1=2l,\qquad l_2=3l,
\tag{14}
\]

上限 \(l_1-h=l\)。为把定积分的来龙去脉写清楚，令

\[
x=3l\sinh^2y.
\]

于是

\[
dx=6l\sinh y\cosh y\,dy,
\qquad
\sqrt{\frac{3l+x}{x}}=\coth y.
\]

当 \(x=0\) 时 \(y=0\)；当 \(x=l\) 时

\[
\sinh y=\frac1{\sqrt3},
\qquad
y=\operatorname{arsinh}\frac1{\sqrt3}
=\frac12\ln3.
\]

因此

\[
\begin{aligned}
\int_0^l\sqrt{\frac{3l+x}{x}}\,dx
&=6l\int_0^{(\ln3)/2}\cosh^2y\,dy\\
&=3l\left[y+\frac12\sinh2y\right]_0^{(\ln3)/2}\\
&=2l\left(1+\frac34\ln3\right).
\end{aligned}
\]

代回式（13），得到

\[
\boxed{
T=
\left(1+\frac34\ln3\right)
\sqrt{\frac{2l}{3a_0}}
}.
\tag{15}
\]

由 \(h=uT=l\)：

\[
\boxed{
u=
\frac{\sqrt{3a_0l/2}}
{1+\frac34\ln3}
}.
\tag{16}
\]

\needspace{3\baselineskip}
\textbf{发射时间的一般闭式与解的唯一性}

令 \(L=l_1-h>0\)，式（13）的定积分可以完整写成

\[
T(L)=\frac{\sqrt{L(L+l_2)}
+l_2\operatorname{arsinh}\sqrt{L/l_2}}
{\sqrt{2a_0l_2}}.
\]

下限的反双曲正弦为零，因此没有遗留的积分常数。以 \(L=l,l_2=3l\)
代入，得到
\(\operatorname{arsinh}(1/\sqrt3)=\tfrac12\ln3\)，恢复式（15）。

对给定 \(u>0\)，未知的 \(h\) 由

\[
H(h)\equiv h-uT(l_1-h)=0
\]

确定。\(H(0)<0\)，\(H(l_1)=l_1>0\)，且 \(H'(h)=1+uT'(l_1-h)>0\)，所以
\(0<h<l_1\) 内只有一个解。 不能把手柄位移和速度作为彼此独立的任选数据。

\needspace{3\baselineskip}
\textbf{手柄外功的能量账本}

随 \(B\) 的惯性系中， \(W_{\rm rel}=\int PS\,dx=\Delta K_{\rm rel}\)。
在炮膛系中 \(dx_A=dx+u\,dt\)，所以气体对 \(A\) 做功为

\[
W_A=W_{\rm rel}+u\int PS\,dt
=W_{\rm rel}+mu\,v_{\rm rel,f}.
\]

它等于 \(\tfrac12m[(u+v_{\rm rel,f})^2-u^2]\)，与式（11）一致。
手柄还要克服气体对 \(B\) 的反向压力，外功为
\(u\int PS\,dt\)。因此不能把炮弹与气体整体当作无外功的封闭系统。
图4-2所涂的曲线下面积是相对功 \(W_{\rm rel}\)，并明确横轴从零起算。

\needspace{3\baselineskip}
\textbf{补充推导与适用条件}

\textbf{为什么可以在 \(B\) 系直接用动能定理。} 手柄保证 \(B\) 始终以恒速
\(u\) 运动，因此随 \(B\) 的参考系仍是惯性系，不需要加入惯性力。气体对
\(A\) 的功只取决于相对位移 \(x\)，而状态 2 时 \(A\)、\(B\)
同速，所以相对初动能为零。若 \(B\)
加速，这两条都不成立，正文的简化不能照搬。

\textbf{绝热做功的上限与量纲。} 对 \(\gamma=2\)，从初始长度 \(l_2\)
膨胀到无限长时，气体最多能做功

\[
W_\infty=P_0Sl_2=ma_0l_2.
\]

式（9）在 \(l_1-h\to\infty\) 时正好趋于这个值。因此相对速度满足

\[
v_{\rm rel}^2<2a_0l_2,
\]

不会因炮管无限长而无限增大。根号内 \(a_0(l_1-h)l_2/(l_1-h+l_2)\)
的量纲为速度平方，也与此上限一致。

\textbf{几何自洽条件。} 必须有

\[
0\le h<l_1,
\]

实际加速过程始终有 \(v_A>u\)，所以出口时相对位移
\(l_1-h>0\)；这一不等式是动力学和出口几何共同给出的可行条件。第三问给出的
\(h=l\)、\(l_1=2l\) 满足该条件。

\textbf{一般绝热指数。} 若 \(\gamma\ne1\)，则

\[
P(x)=P_0\left(\frac{l_2}{l_2+x}\right)^\gamma,
\]

相对动能由 \(\int_0^xP(x')S\,dx'\) 给出。\(\gamma=2\)
的特殊价值在于积分成为简单有理函数，并使后续时间积分能够化成题给的根式公式。

\needspace{3\baselineskip}
\textbf{任意手柄速度时的唯一解与计算方法}

令 \(L=l_1-h\)。从前面的相对速度式可写出

\[
v_{\rm rel}(L)=\sqrt{\frac{2a_0Ll_2}{L+l_2}},\qquad
T(L)=\int_0^L\frac{dx}{v_{\rm rel}(x)}.
\]

实际手柄匀速移动距离 \(h=uT(L)\)。定义 \(H(h)=h-uT(l_1-h)\)，在
\(0<h<l_1\) 中有

\[
H'(h)=1+\frac{u}{v_{\rm rel}(l_1-h)}>0.
\]

而 \(H(0)=-uT(l_1)<0\)、\(H(l_1)=l_1>0\)，所以对每个 \(u>0\)
都有唯一物理解。这个证明排除了迭代时误选根的可能，也说明图4-2中相对位移区间必须以
\(L\) 为上端，不能把手柄移动的 \(h\) 再加进气体膨胀长度。

可在 \([0,l_1]\) 二分求根。Newton法的更新式为

\[
h_{j+1}=h_j-
\frac{h_j-uT(l_1-h_j)}{1+u/v_{\rm rel}(l_1-h_j)}.
\]

迭代点应保持在物理区间；靠近 \(h=l_1\)
时相对末速度变小，采用二分夹逼更容易控制误差。得到 \(h\) 后，先求 \(L\)
和相对末速度，再恢复地面速度 \(v_A=u+v_{\rm rel}\)。

两个极限也可复核方向。\(u\to0\) 时
\(h\to0,L\to l_1\)，回到后弹固定的气体膨胀问题。\(u\to\infty\) 时
\(L\ll l_2\)，有 \(T(L)\simeq\sqrt{2L/a_0}\)，因而

\[
L\simeq\frac{a_0l_1^2}{2u^2},\qquad
v_A\simeq u+\frac{a_0l_1}{u}.
\]

大手柄速度会使前弹很快到炮口，留给相对膨胀的时间很短；前弹在地面系仍继承手柄所带来的运动。气体对相对运动所做的功和外部手柄对装置所做的功应分开列账，不能以``气体储能固定''为由否定这个速度极限。

\subsection{六、关键洞察}\label{ux516dux5173ux952eux6d1eux5bdf-3}

第二问最关键的是使用随 \(B\)
匀速运动的惯性系：初始相对速度为零，气柱长度只增加
\(l_1-h\)，复杂的移动边界立即变成固定起点的做功问题。

\subsection{七、易错点}\label{ux4e03ux6613ux9519ux70b9-3}

\begin{enumerate}
\def\labelenumi{\arabic{enumi}.}
\tightlist
\item
  把绝热关系写成 \(P(l_2+x)=\text{常量}\)，漏掉 \(\gamma=2\)。
\item
  压强功漏乘截面积 \(S\)。
\item
  第二问仍把气柱伸长写成 \(l_1\)。
\item
  求得相对速度后忘记加回 \(u\)。
\item
  时间积分上限写成 \(l_1\) 而非 \(l_1-h\)。
\item
  第一问把 \(A\) 在状态 2 的初速度误当成零；只有切换到 \(B\)
  系后相对初速度才为零。
\item
  把压强 \(P_0\) 误当作大气压；题中 \(P_0=ma_0/S\) 是状态 2 的气体压强。
\item
  用气体对两颗炮弹的总功直接等于 \(A\) 的动能增量，没有处理手柄对 \(B\)
  的外功。
\item
  在根式中允许 \(h\ge l_1\)，忽略两炮弹相遇后模型已经失效。
\item
  求出 \(T\) 后忘记再用 \(h=uT\)，把中间量当成第三问最终答案。
\end{enumerate}

\subsection{八、课堂讲解建议}\label{ux516bux8bfeux5802ux8bb2ux89e3ux5efaux8bae-3}

先用图 4-1 逐状态口述：状态 2 时两弹同速 \(u\)，间距 \(l_2\)；状态 3 时
\(A\) 前进 \(l_1\) 到炮口，\(B\) 前进 \(h\)，所以间距只增加
\(l_1-h\)。要求学生在任何公式之前写出这一几何关系。

接着在 \(P\)---\(x\) 图上画
\(P=P_0l_2^2/(l_2+x)^2\)。用曲线下面积解释压强功，并先估计面积小于矩形
\(P_0Sx\)，从而给结果设置上界。第一问留在炮膛系，第二问切换到匀速 \(B\)
系；将两种参考系的初速度条件并排比较。

第三问先从动能定理得到 \(v_{\rm rel}(x)\)，再用
\(dt=dx/v_{\rm rel}\)。积分下限附近 \(v_{\rm rel}\propto\sqrt x\)，所以
\(dt\propto dx/\sqrt x\)，积分虽有形式奇点但仍收敛。完成双曲代换后，用
\(h=uT\) 闭合，而不是把 \(h\)、\(u\) 当成独立量。

最后检查三项：\(u\) 的量纲是速度；\(u>0\)；所得
\(h=l<l_1=2l\)。三项都满足才算完成。

\subsection{九、变式训练}\label{ux4e5dux53d8ux5f0fux8badux7ec3-3}

将绝热指数改为任意 \(\gamma\ne1\)。压强功为

\[
\begin{aligned}
W(x)
&=P_0Sl_2^\gamma
\int_0^x(l_2+x')^{-\gamma}\,dx'\\
&=\frac{P_0Sl_2}{\gamma-1}
\left[1-\left(\frac{l_2}{l_2+x}\right)^{\gamma-1}\right].
\end{aligned}
\]

利用 \(P_0S=ma_0\)：

\[
\boxed{
v_{\rm rel}^2(x)
=\frac{2a_0l_2}{\gamma-1}
\left[1-\left(\frac{l_2}{l_2+x}\right)^{\gamma-1}\right]
}.
\]

令 \(\gamma=2\) 恢复正文式（12）。若手柄速度随时间变化，\(B\)
系不再是惯性系，需要加入 \(-m\dot u\)
的惯性力；也可以始终留在炮膛系，同时把手柄对 \(B\)
的功纳入系统能量收支。

在一般 \(\gamma\) 下，从相对位移 \(0\) 运动到 \(L=l_1-h\) 所需的时间为

\[
T_\gamma
=\int_0^L
\frac{dx}{
\sqrt{\dfrac{2a_0l_2}{\gamma-1}
\left[1-\left(\dfrac{l_2}{l_2+x}\right)^{\gamma-1}\right]}}.
\]

令 \(\xi=x/l_2\)、\(\Lambda=L/l_2\)，可把尺度和形状完全分开：

\[
T_\gamma
=\sqrt{\frac{(\gamma-1)l_2}{2a_0}}
\int_0^\Lambda
\frac{d\xi}{
\sqrt{1-(1+\xi)^{1-\gamma}}}.
\]

根号前的因子具有时间量纲，积分只取决于无量纲比值 \(\Lambda\) 和绝热指数
\(\gamma\)。当 \(\xi\to0\) 时，

\[
1-(1+\xi)^{1-\gamma}
=(\gamma-1)\xi+O(\xi^2),
\]

因而被积函数只按 \(1/\sqrt\xi\) 发散，积分仍然收敛。原题取
\(\gamma=2\)时，这个无量纲积分恰可用双曲代换化为对数闭式；其他
\(\gamma\) 通常保留为积分或数值求值。最后仍要与 \(h=uT_\gamma\)
联立，因为 \(L=l_1-h\) 也含未知量 \(h\)。

\newpage

\section{第 5 题}\label{ux7b2c-5-ux9898}

已知四边形 \(ABCD\) 是一个正方形物体通过理想薄凸透镜所成的像。
请根据该四边形，用尺规作出原正方形和透镜的位置，并写出作图的关键步骤。

\begin{figure}[H]
\centering
\includegraphics[width=0.767\linewidth,height=\textheight,keepaspectratio,alt={图5-1　由同一薄透镜公式生成的正方形像四边形}]{../figures/question-05.pdf}
\caption*{图5-1　由同一薄透镜公式生成的正方形像四边形}
\end{figure}

\subsection{一、题型定位}\label{ux4e00ux9898ux578bux5b9aux4f4d-4}

两组平行边的像线先给出灭点连线。邻边和对角线的两组正交条件确定光心候选；再由焦平面迹线作透镜与焦点，用两条主光线逐点恢复原正方形。

\subsection{二、知识点}\label{ux4e8cux77e5ux8bc6ux70b9-4}

\begin{enumerate}
\def\labelenumi{\arabic{enumi}.}
\tightlist
\item
  过薄透镜光心的光线不偏折。
\item
  平行物方直线的像线交于该方向的像方焦平面点。
\item
  原正方形邻边及两条对角线分别互相垂直。
\item
  以线段为直径的圆上的角为直角。
\item
  物方平行主轴的光线经透镜后通过像方焦点。
\end{enumerate}

\subsection{三、物理图像}\label{ux4e09ux7269ux7406ux56feux50cf-4}

这里把原正方形、像四边形和主光路放在同一个作图平面中，光心也在该
平面内。二维图中的透镜是一条通过光心的直线。若另把四边形理解成
三维像平面的投影，光心未必在给出的平面中，就不能直接把空间直角
写成纸面上的直角圆；需要额外的空间标定。本解采用原题参考作图的平面模型。

\subsection{四、解题入口}\label{ux56dbux89e3ux9898ux5165ux53e3-4}

先作 \(P=AB\cap CD\)、\(Q=AD\cap BC\)，连接 \(\ell=PQ\)。
对角线的灭点可直接取 \(R=AC\cap\ell\)、\(S=BD\cap\ell\)，
不必另猜中点。分别以 \(PQ,RS\) 为直径作圆，取公共点作为光心候选。

\subsection{五、解析}\label{ux4e94ux89e3ux6790-4}

\begin{figure}[H]
\centering
\includegraphics[width=0.85\linewidth,height=\textheight,keepaspectratio,alt={图5-2　灭点、两直角圆、光心及物像主光线的完整恢复}]{../figures/solution-05.pdf}
\caption*{图5-2　灭点、两直角圆、光心及物像主光线的完整恢复}
\end{figure}

\needspace{3\baselineskip}
\textbf{（一）求边方向和对角线方向的灭点}

原正方形的两组对边分别平行。它们的像线相交于

\[
P=AB\cap CD,\qquad Q=AD\cap BC.
\tag{1}
\]

\(OP,OQ\) 分别平行于两组物方边方向，所以

\[
OP\perp OQ.
\tag{2}
\]

所有物方无穷远方向的像，都落在像方焦平面迹线 \(\ell\) 上； 两点 \(P,Q\)
确定它。两条物方对角线也互相垂直，其灭点为

\[
R=AC\cap\ell,\qquad S=BD\cap\ell,
\tag{3}
\]

并满足 \(OR\perp OS\)。作以 \(PQ\)、\(RS\) 为直径的圆
\(\Gamma_1,\Gamma_2\)。由泰勒斯定理，

\[
\boxed{O\in\Gamma_1\cap\Gamma_2.}
\tag{4}
\]

一般有两个候选光心，它们关于 \(\ell\) 对称，先全部保留。

\needspace{3\baselineskip}
\textbf{（二）由焦平面迹线确定透镜和焦点}

对每个候选 \(O\)，过 \(O\) 作平行于 \(\ell\) 的直线 \(L\)，
它就是透镜在作图平面中的迹线。过 \(O\) 作垂直 \(\ell\) 的主轴， 与
\(\ell\) 交于像方焦点 \(F\)。将 \(F\) 关于 \(O\) 对称得到物方焦点
\(F'\)。

\[
\boxed{L\parallel PQ,\quad O\in L,\quad
F=\ell\cap(\text{过}O\text{的}\ell\text{垂线}),\quad f=OF.}
\tag{5}
\]

焦距在这个平面模型中由已给图形确定，不能再任意选物距或正方形尺度。
若改变物方尺度却保持光心和像四边形不动，通常会破坏同一透镜的 成像关系。

\needspace{3\baselineskip}
\textbf{（三）逐点恢复原正方形}

对像点 \(A\)：

\begin{enumerate}
\def\labelenumi{\arabic{enumi}.}
\tightlist
\item
  作直线 \(FA\)，与透镜迹线 \(L\) 交于 \(H_A\)。
\item
  过 \(H_A\) 作平行主轴的直线。
\item
  它与 \(OA\) 的交点为原物点 \(A'\)。
\end{enumerate}

这对应``物方平行主轴的光线，经透镜后通过像方焦点''和``过光心
不偏折的光线''两条主光线。对 \(B,C,D\) 作同样操作，得到 \(B',C',D'\)。

\[
\boxed{
H_X=FX\cap L,\qquad
X'=OX\cap(\text{过}H_X\text{的主轴平行线}),
\quad X=A,B,C,D.}
\tag{6}
\]

最后检查 \(A'B'C'D'\) 的邻边垂直、相邻边等长，以及四点与各自像点、
光心共线。根据实像位于物体另一侧等成像条件选择候选构型。
如果题目未额外指定透镜所在侧，两个几何构型可能都需要报告。

\needspace{3\baselineskip}
\textbf{（四）用薄透镜公式核验作图}

以 \(O\) 为原点，取主轴单位矢量 \(\boldsymbol n\) 指向 \(\ell\)， 令像点
\(X\) 的有向轴坐标为 \(v=\boldsymbol n\cdot X\)。 像方焦平面为
\(v=f\)。上面的作图给出

\[
\boxed{X'=-\frac{f}{v-f}X.}
\tag{7}
\]

当 \(v>f>0\) 时，物点在相反一侧，其正物距为

\[
u=-\boldsymbol n\cdot X'=\frac{fv}{v-f}.
\]

于是

\[
\frac1u+\frac1v
=\frac{v-f}{fv}+\frac1v=\frac1f.
\tag{8}
\]

横向坐标也按 \(-u/v\) 缩放，符合放大率。式（7）可逐点验证尺规
结果，而不是只凭四条射线看起来像正方形就结束。

\needspace{3\baselineskip}
\textbf{（五）为什么恢复后是正方形}

把像点沿 \(OX\) 逆变换到物方后，像线 \(AB,CD\) 的共同灭点 \(P\)
退到同一物方无穷远方向，故 \(A'B'\parallel C'D'\parallel OP\)；
另一组同理平行于 \(OQ\)。由第一直角圆，两方向垂直，物方四边形是矩形。

第二直角圆保证两物方对角线垂直。若矩形的邻边向量为
\(\boldsymbol a,\boldsymbol b\)，则两对角线向量为
\(\boldsymbol a+\boldsymbol b\)、\(\boldsymbol a-\boldsymbol b\)，其点积为

\[
(\boldsymbol a+\boldsymbol b)\cdot(\boldsymbol a-\boldsymbol b)
=|\boldsymbol a|^2-|\boldsymbol b|^2.
\]

点积为零给 \(|\boldsymbol a|=|\boldsymbol b|\)，所以恢复后确为正方形。
边的正交和对角线正交各有作用，不能漏掉后一条件。

\needspace{3\baselineskip}
\textbf{（六）调和关系的辅助检查}

设 \(E=AC\cap BD\)。物方 \(E'\) 是两对角线的中点；中点与该直线无穷远点
构成调和点列。中心成像保持交比，所以

\[
(A,C;E,R)=-1,\qquad(B,D;E,S)=-1.
\]

因此 \(R,S\) 也可由调和共轭作出，但它们已经由式（3）确定。 若在 \(AC\)
上取仿射坐标 \(A=-1,C=1,E=e\)，调和点坐标为 \(R=1/e\)。 当 \(E=0\)
时灭点退到无穷远；当 \(e\) 接近端点 \(\pm1\) 时，\(1/e\)
也接近同一端点，不能说它必向另一侧大幅移动。

\needspace{3\baselineskip}
\textbf{一组可复算的坐标示例}

图5-1的像由焦距 \(f=1\)、透镜线 \(x=0\)、光心 \((0,0)\) 生成。
物方正方形中心为 \((-2.6,0.1)\)，半边长为1，边方向角为
\(0.5\,\mathrm{rad}\)。记

\[
\boldsymbol t=(\cos0.5,\sin0.5),\qquad
\boldsymbol s=(-\sin0.5,\cos0.5).
\]

物方四点按顺序取 \(A'=C_0-\boldsymbol t-\boldsymbol s\)、
\(B'=C_0+\boldsymbol t-\boldsymbol s\)、
\(C'=C_0+\boldsymbol t+\boldsymbol s\)、
\(D'=C_0-\boldsymbol t+\boldsymbol s\)。 将物方轴坐标记为
\(x'<-1\)，像点由

\[
(x,y)=\frac{1}{x'+1}(x',y')
\]

生成。这样，图中的顶点、灭点、直角圆和恢复物点属于同一成像系统。
其焦平面迹线为 \(x=1\)，四个灭点均在这条线上，且
\(\boldsymbol{OP}\cdot\boldsymbol{OQ}=0\)、
\(\boldsymbol{OR}\cdot\boldsymbol{OS}=0\)。

\needspace{3\baselineskip}
\textbf{用实际四点回查透镜恢复}

图5-2的数据取透镜光心 \(O=(0,0)\)、焦距 \(f=1\)，物方正方形中心
\((-2.6,0.1)\)，半边长 \(1\)，与轴线夹角 \(0.5\) 弧度。取正交边方向
\(\boldsymbol t=(\cos0.5,\sin0.5)\)、\(\boldsymbol s=(-\sin0.5,\cos0.5)\)，四个物点为中心加减这两个向量，故边长和邻边垂直可在物方直接核验。

对每一点用实验室坐标映射
\(X'=X/(X+1),Y'=Y/(X+1)\)，得到图中的四个像点。它们不是把原正方形沿某个方向简单压缩形成。两组边的消失点分别为

\[
P=(1,\tan0.5),\qquad Q=(1,-\cot0.5),
\]

两组对角线的消失点分别为

\[
R=(1,\tan(0.5+\pi/4)),\qquad
S=(1,\tan(0.5-\pi/4)).
\]

因此
\(OP\cdot OQ=0\)、\(OR\cdot OS=0\)，光心确在两个直角圆上；四个点都在像方焦平面
\(X'=1\)。按这两个圆作交点会得到 \(O=(0,0)\) 和另一候选
\((2,0)\)，还需依真实光线方向与物像性质选支。

取任一像点 \(X'\)，令 \(H=FX'\cap L\)，则 \(H\)
是物点发出的轴平行光在透镜上的入射点。用光心直线 \(OX'\) 与过 \(H\)
的轴平行线相交恢复物点。将此作图逐个用于A、B、C、D，重新得到原正方形。这个回代同时检查了焦距、透镜线方向、两个焦点和正方形的尺度。

一般射影四边形不能保证满足两个直角圆的共同光学条件，图中的辅助圆也不是任意选点连成。原稿只给形状后再声称尺度任意的做法已撤去；本题二维理想薄透镜构型的长度由焦距与已知像点坐标确定。

\needspace{3\baselineskip}
\textbf{作图的退化与测量误差}

消失点采用延长线交点。若某组像边接近平行，该消失点会很远；此时应把它理解为射影平面中的无穷远点，原来的有限直径圆作图不能直接照搬。数值处理可用齐次坐标：把点写为
\((X,Y,1)\)，两点叉乘给直线，两条直线叉乘给交点。交点最后一个分量接近零时，先记录方向，避免除以一个很小的数。

两个直角圆还可能相切或重合。一般构型有两个光心候选；圆相切时两个候选合并；圆重合时单凭这两个直角约束不足以确定光心。题目若恰处于这些退化位置，需要增加已知焦距、实际透镜方向或另一组独立物像信息，不能报告唯一的透镜位置。

对图5-2的非退化例，选出光心后，焦距为光心到消失线的垂距；焦平面法线决定主轴。因此在尺度已经标定的像平面内，焦距也有确定长度。若只是把没有尺标的图片整体放大，所有几何长度一起按比例放大，那是输入数据的尺度改变，不是固定像点条件下可以任意改变透镜焦距。

实验作图中四个像点有测量误差，两圆约束往往不能精确交会。可先由两组边估计消失线，再用正方形的邻边垂直、对角线垂直和逆映射后四边等长共同检验候选解。越接近退化位置，延长线交点对端点误差越敏感；此时多画小数位并不能提高物理精度。正文尺规构造针对理想无误差数据，图中数值回代是对这套模型的检查。

\subsection{六、关键洞察}\label{ux516dux5173ux952eux6d1eux5bdf-4}

两条直角圆确定光心候选，灭点连线确定焦平面迹线。
后者又确定透镜迹线及焦距，因此恢复物点还必须使用折射主光线，
不能在光心射线上任意选一个相似正方形。

\subsection{七、易错点}\label{ux4e03ux6613ux9519ux70b9-4}

\begin{enumerate}
\def\labelenumi{\arabic{enumi}.}
\tightlist
\item
  把对角线交点当成光心。
\item
  只有一个直角圆便声称光心唯一。
\item
  将纸面上的直角圆无条件用于三维投影问题。
\item
  恢复物点时任意选尺度，未检查同一焦距。
\item
  忽略 \(PQ\) 是像方焦平面迹线，漏掉透镜位置。
\item
  两个候选光心中只画一个，不说明物像侧条件。
\item
  对平行像边强行求有限灭点，应改用无穷远方向的退化构造。
\end{enumerate}

\subsection{八、课堂讲解建议}\label{ux516bux8bfeux5802ux8bb2ux89e3ux5efaux8bae-4}

板书按``灭点---光心---透镜及焦点---物点''四步组织。
第一圆恢复邻边直角，第二圆恢复等边；求得光心后，要继续由 \(PQ\)
作透镜平行线和主轴垂线。

逐点恢复时给 \(OA\) 和 \(FH_A\) 两条光线不同颜色，请学生说明各自
使用的成像规律。最后用式（7）、（8）检查同一透镜公式是否对四个点
都成立。原稿图中任意摆放的像和正方形无法提供这种检查，现图均由
同一坐标模型生成。

对于空间模型的讨论，若像平面与光心相距 \(d\)，在像平面中画出的 \(P,Q\)
的二维位移向量为 \(\boldsymbol p,\boldsymbol q\)，空间正交条件是
\(\boldsymbol p\cdot\boldsymbol q+d^2=0\)，而非二维点积为零。
这解释了为什么必须先声明作图平面。

\subsection{九、变式训练}\label{ux4e5dux53d8ux5f0fux8badux7ec3-4}

若原物是长方形，只能直接使用边方向的直角圆；对角线通常不垂直。
已知长宽比时可把它转为对角线夹角的度量约束。若原物为菱形，则保留
对角线正交而不能保留邻边正交。

若 \(AB\parallel CD\)，\(P\) 在无穷远，\(OP\) 表示已知平行方向；
\(OP\perp OQ\) 变成过 \(Q\) 的一条垂线。焦平面迹线也应与该方向平行。
将这条垂线与另一有限直角圆联立，再按焦平面、透镜、主光线的顺序恢复。

若两直角圆没有实交点，说明给出的精确平面四边形不符合本模型中的
正方形薄透镜成像。若有两个候选，分别恢复并检查物像侧与焦距正号；
候选筛选是作图的一部分。

\newpage

\section{第 6 题}\label{ux7b2c-6-ux9898}

一根长为 \(l\)、质量为 \(m\) 的均匀刚性细杆在桌边一个很细的固定凸起上方
\(h\)
处保持水平，并由静止释放。细杆与凸起发生完全非弹性碰撞，碰撞点为细杆右侧三等分点。

\begin{enumerate}
\def\labelenumi{\arabic{enumi}.}
\tightlist
\item
  求碰撞后瞬间细杆质心速度与绕质心转动的角速度。
\item
  杆与凸起间的摩擦因数为
  \(\mu=2\)。若碰撞后瞬间接触点不相对滑动，求允许的最大释放高度
  \(h_{\max}\)。
\item
  取
  \(h=h_{\max}/2\)，求碰撞后杆与凸起刚要发生相对滑动时，杆与水平方向的夹角
  \(\theta_0\)。
\item
  写出发生相对滑动后细杆的运动微分方程。
\end{enumerate}

\begin{figure}[H]
\centering
\includegraphics[width=0.621\linewidth,height=\textheight,keepaspectratio,alt={图6-1　水平细杆自由下落并与桌边凸起碰撞}]{../figures/question-06.pdf}
\caption*{图6-1　水平细杆自由下落并与桌边凸起碰撞}
\end{figure}

\subsection{一、题型定位}\label{ux4e00ux9898ux578bux5b9aux4f4d-5}

细杆先自由下落，与偏离质心的凸起碰撞后开始转动。碰撞阶段用对凸起的角动量守恒；随后分别反求法向力和静摩擦，找到失去固定接触的临界角。

\subsection{二、知识点}\label{ux4e8cux77e5ux8bc6ux70b9-5}

\begin{enumerate}
\def\labelenumi{\arabic{enumi}.}
\tightlist
\item
  碰撞极短时间内对碰撞点取角动量守恒。
\item
  完全非弹性碰撞后的瞬间，接触点法向相对速度为零；若不滑动，则接触点瞬时速度为零。
\item
  固定接触阶段，细杆绕凸起作定轴转动。
\item
  静摩擦大小由动力学决定，并满足 \(|f|\le \mu N\)。
\item
  滑动后应使用质心平动方程与绕质心转动方程联立。
\end{enumerate}

\subsection{三、物理图像}\label{ux4e09ux7269ux7406ux56feux50cf-5}

碰撞点位于右侧三等分点，即它到右端的距离为
\(l/3\)。细杆质心在中点，因此碰撞点 \(O\) 到质心 \(C\) 的距离为

\[
d=\frac l2-\frac l3=\frac l6.
\tag{1}
\]

碰撞后若接触点粘住，质心沿以 \(O\) 为圆心、半径 \(d\)
的圆周运动；随着细杆转动，保持这种运动所需的切向与法向接触力不断变化，最终可能达到库仑摩擦极限。

\subsection{四、解题入口}\label{ux56dbux89e3ux9898ux5165ux53e3-5}

碰撞阶段对凸起 \(O\) 取矩，冲量不产生角冲量；碰撞后的固定阶段对 \(O\)
写转动方程，同时把质心加速度分解为径向与切向分量，反求
\(f,N\)。这样摩擦临界条件中的每一个力都有明确来源。

\subsection{五、解析}\label{ux4e94ux89e3ux6790-5}

\begin{figure}[H]
\centering
\includegraphics[width=0.86\linewidth,height=\textheight,keepaspectratio,alt={图6-2　固定接触阶段的几何关系与受力分解}]{../figures/solution-06.pdf}
\caption*{图6-2　固定接触阶段的几何关系与受力分解}
\end{figure}

\needspace{3\baselineskip}
\textbf{（一）碰撞后瞬间的速度}

碰撞前细杆自由落体，由机械能守恒

\[
mgh=\frac12mv_0^2,
\tag{2}
\]

故

\[
v_0=\sqrt{2gh}.
\tag{3}
\]

均匀细杆绕质心的转动惯量为

\[
I_C=\frac1{12}ml^2.
\tag{4}
\]

绕接触点 \(O\) 的转动惯量由平行轴定理得到

\[
\begin{aligned}
I_O
&=I_C+md^2\\
&=\frac1{12}ml^2+m\left(\frac l6\right)^2\\
&=\left(\frac3{36}+\frac1{36}\right)ml^2\\
&=\frac19ml^2.
\end{aligned}
\tag{5}
\]

碰撞时间极短，重力冲量可忽略。凸起冲量通过 \(O\)，对 \(O\)
的冲量矩为零。因此碰撞前后对 \(O\) 角动量守恒：

\[
mv_0d=I_O\omega_1.
\tag{6}
\]

代入 \(d=l/6\) 与式（5）：

\[
\begin{aligned}
\omega_1
&=\frac{mv_0(l/6)}{ml^2/9}\\
&=\frac{3v_0}{2l}\\
&=\boxed{\frac{3\sqrt{2gh}}{2l}}.
\end{aligned}
\tag{7}
\]

碰撞后若接触点瞬时不滑动，杆绕 \(O\) 转动，所以质心速度

\[
\begin{aligned}
V_C
&=\omega_1d\\
&=\frac{3v_0}{2l}\cdot\frac l6\\
&=\boxed{\frac{v_0}{4}=\frac{\sqrt{2gh}}4}.
\end{aligned}
\tag{8}
\]

这里的 \(V_C\) 方向垂直于
\(OC\)，在碰撞后瞬间即为竖直方向。还可以用动能损失检查``完全非弹性''是否被正确处理。碰撞前动能为

\[
K_-=\frac12mv_0^2,
\]

碰撞后绕 \(O\) 瞬时转动的动能为

\[
\begin{aligned}
K_+
&=\frac12I_O\omega_1^2\\
&=\frac12\cdot\frac19ml^2
\left(\frac{3v_0}{2l}\right)^2\\
&=\frac18mv_0^2.
\end{aligned}
\]

因此

\[
\Delta K=K_--K_+=\frac38mv_0^2>0.
\]

角动量守恒而机械能显著减少，正好符合完全非弹性碰撞的物理特征；若算出
\(K_+=K_-\)，说明误把碰撞过程当成了弹性过程。

\needspace{3\baselineskip}
\textbf{（二）碰撞后瞬间不滑动的高度条件}

碰撞后紧接着杆仍水平。以质心指向凸起的方向为杆轴方向，此时质心的向心加速度大小为
\(\omega_1^2d\)。沿水平方向只有摩擦力，因此

\[
f=m\omega_1^2d.
\tag{9}
\]

代入式（7）：

\[
\begin{aligned}
f
&=m\left(\frac{3\sqrt{2gh}}{2l}\right)^2\frac l6\\
&=m\frac{18gh}{4l^2}\frac l6\\
&=\frac{3mgh}{4l}.
\end{aligned}
\tag{10}
\]

重力对 \(O\) 的力矩大小为 \(mgd\)，故碰撞后瞬间的角加速度

\[
\begin{aligned}
\beta
&=\frac{mgd}{I_O}\\
&=\frac{mg(l/6)}{ml^2/9}\\
&=\frac{3g}{2l}.
\end{aligned}
\tag{11}
\]

质心竖直加速度大小为切向加速度 \(\beta d\)：

\[
\beta d=\frac{3g}{2l}\cdot\frac l6=\frac g4.
\tag{12}
\]

质心向下加速，所以竖直方向

\[
mg-N=m\frac g4,
\tag{13}
\]

从而

\[
N=\frac34mg.
\tag{14}
\]

静摩擦可维持接触点不滑动的条件为

\[
f\le \mu N.
\tag{15}
\]

将式（10）、（14）与 \(\mu=2\) 代入：

\[
\frac{3mgh}{4l}\le2\cdot\frac34mg.
\tag{16}
\]

消去公共因子 \(3mg/4\)：

\[
\frac hl\le2.
\tag{17}
\]

故

\[
\boxed{h_{\max}=2l}.
\tag{18}
\]

\needspace{3\baselineskip}
\textbf{（三）固定接触阶段的临界角}

题设取

\[
h=\frac12h_{\max}=l.
\tag{19}
\]

设细杆从水平位置逆时针转过
\(\theta\)，左端下降、右端上升。质心下降高度为
\(d\sin\theta\)。固定接触阶段机械能守恒：

\[
\frac12I_O\omega^2
=\frac12I_O\omega_1^2+mgd\sin\theta.
\tag{20}
\]

两边乘 \(2/I_O\)，代入 \(I_O=ml^2/9\)、\(d=l/6\) 以及 \(h=l\) 时

\[
\omega_1^2=\frac{9g}{2l},
\tag{21}
\]

得到

\[
\begin{aligned}
\omega^2
&=\frac{9g}{2l}
+\frac{2mg(l/6)}{ml^2/9}\sin\theta\\
&=\frac gl\left(\frac92+3\sin\theta\right).
\end{aligned}
\tag{22}
\]

沿杆轴、指向凸起的方向列质心径向方程：

\[
f-mg\sin\theta=m\omega^2d.
\tag{23}
\]

代入式（22）：

\[
\begin{aligned}
f
&=mg\sin\theta
+m\frac gl\left(\frac92+3\sin\theta\right)\frac l6\\
&=mg\left(\frac34+\frac32\sin\theta\right).
\end{aligned}
\tag{24}
\]

重力对 \(O\) 的力矩为 \(mgd\cos\theta\)，所以

\[
\beta=\frac{mgd\cos\theta}{I_O}
=\frac{3g}{2l}\cos\theta.
\tag{25}
\]

沿杆法向列质心切向方程：

\[
mg\cos\theta-N=m\beta d.
\tag{26}
\]

代入式（25）：

\[
\begin{aligned}
N
&=mg\cos\theta
-m\frac{3g}{2l}\cos\theta\frac l6\\
&=\frac34mg\cos\theta.
\end{aligned}
\tag{27}
\]

刚要滑动时 \(f=\mu N=2N\)。由式（24）、（27）：

\[
\frac34+\frac32\sin\theta_0
=\frac32\cos\theta_0.
\tag{28}
\]

同除以 \(3/4\)：

\[
1+2\sin\theta_0=2\cos\theta_0.
\tag{29}
\]

令 \(t=\tan(\theta_0/2)\)，使用

\[
\sin\theta_0=\frac{2t}{1+t^2},
\qquad
\cos\theta_0=\frac{1-t^2}{1+t^2},
\tag{30}
\]

代入式（29）：

\[
1+t^2+4t=2-2t^2,
\tag{31}
\]

即

\[
3t^2+4t-1=0.
\tag{32}
\]

取正根

\[
t=\frac{\sqrt7-2}{3}.
\tag{33}
\]

因此

\[
\boxed{
\theta_0=2\arctan\frac{\sqrt7-2}{3}
\approx24.3^\circ
}.
\tag{34}
\]

\needspace{3\baselineskip}
\textbf{（四）滑动后的运动微分方程}

滑动后，凸起不再是杆的固定转轴。取 \(x>0\) 为从凸起 \(O\) 指向左侧质心
\(C\) 的距离，\(\theta\) 为杆右向轴从水平逆时针转过的角度。于是
\(\overrightarrow{OC}=(-x\cos\theta,-x\sin\theta)\)。质心在随杆转动的基底中的加速度分量为

\[
a_{\parallel}=\ddot x-x\dot\theta^2,
\qquad
a_{\perp}=x\ddot\theta+2\dot x\dot\theta.
\tag{35}
\]

设 \(N\) 为法向力，滑动摩擦力大小
\(f=\mu N\)，方向与接触点沿杆相对滑动方向相反。若当前滑动方向对应摩擦沿
\(-x\) 方向，则

\[
m(\ddot x-x\dot\theta^2)=mg\sin\theta-\mu N.
\tag{36}
\]

杆法向方向有

\[
m(x\ddot\theta+2\dot x\dot\theta)=mg\cos\theta-N.
\tag{37}
\]

绕质心列转动方程，接触力的法向分量对质心力矩大小为 \(Nx\)：

\[
I_C\ddot\theta=Nx,
\qquad
I_C=\frac1{12}ml^2.
\tag{38}
\]

式（36）---（38）连同初始条件 \(x=d\)、\(\theta=\theta_0\)
以及临界瞬间的
\(\dot x,\dot\theta\)，构成滑动后的闭合微分方程组。若计算中发现
\(\dot x\) 变号，式（36）中摩擦项的符号也必须随之改变。

\needspace{3\baselineskip}
\textbf{运动方向、滑动符号和初值闭合}

图6-2的质心位于凸起左侧。重力使左端下落，因此 \(\theta\) 应沿逆时针
增加，质心坐标为 \((-d\cos\theta,-d\sin\theta)\)，下降量才是
\(d\sin\theta\)。若图画左端上升，却仍写重力势能减少，会把图和推导割裂。

滑动后取

\[
\boldsymbol e_\parallel=(-\cos\theta,-\sin\theta),\qquad
\boldsymbol e_\perp=(\sin\theta,-\cos\theta).
\]

法向力为 \(-N\boldsymbol e_\perp\)，摩擦沿
\(-\operatorname{sgn}(\dot x)\boldsymbol e_\parallel\)，完整方程为

\[
\begin{cases}
\ddot x-x\dot\theta^2
=g\sin\theta-\mu N\,\operatorname{sgn}(\dot x)/m,\\
x\ddot\theta+2\dot x\dot\theta
=g\cos\theta-N/m,\\
\ddot\theta=Nx/I_C.
\end{cases}
\]

在首次滑动点 \(x=d\)、\(\dot x=0\)，
\(\dot\theta=\sqrt{(g/l)(9/2+3\sin\theta_0)}\)。 向外滑动取 \(\dot x>0\)
分支；由临界静摩擦的继续增长可确定此分支， 不能把
\(\operatorname{sgn}(0)\) 简单设为零后当成无摩擦。

把后两行联立消去角加速度，得到

\[
N=\frac{m(g\cos\theta-2\dot x\dot\theta)}
{1+mx^2/I_C},\qquad
\ddot\theta=\frac{Nx}{I_C}.
\]

再代入第一行就是两个坐标的闭合初值问题。必须同时检查 \(N\ge0\) 和 \(x\)
仍位于杆上；一旦 \(N=0\) 后需要负法向力，改为无接触自由运动。

碰撞瞬间只有竖直法向冲量。由质心速度从 \(-v_0\) 变为 \(-v_0/4\)， 冲量为
\(3mv_0/4\) 向上；绕质心的角冲量恰给式（7）。
水平初始接触速度为零，因此冲量阶段无需水平摩擦冲量。
这也说明第二问的高度条件来自碰后有限时间静摩擦，而不是冲量比。

\needspace{3\baselineskip}
\textbf{补充推导与适用条件}

\textbf{冲量阶段与有限时间动力学必须分开。}
碰撞持续时间趋于零，凸起冲量可以很大，重力冲量却可忽略，所以对 \(O\)
角动量守恒；碰撞结束后，冲量消失，接触力恢复为有限的
\(N,f\)，此时才用牛顿第二定律判断静摩擦是否足够。把式（9）、（13）用于碰撞冲量，会在量纲上把``力''与``冲量''混在一起。

\textbf{一般摩擦因数下的最大高度。} 碰撞后杆仍水平时，正文已经得到

\[
f=\frac{3mgh}{4l},
\qquad
N=\frac34mg.
\]

因此

\[
\frac{f}{N}=\frac hl.
\]

对任意静摩擦因数 \(\mu\)，不立即滑动的条件为 \(h/l\le\mu\)，即

\[
\boxed{h_{\max}=\mu l}.
\]

代入 \(\mu=2\) 才得到本题的 \(2l\)。本题的长度关系由这一摩擦临界式给出。

\textbf{一般临界角方程。} 若释放高度仍记为
\(h\)，固定接触阶段能量守恒给出

\[
\omega^2=\frac{9gh}{2l^2}+\frac{3g}{l}\sin\theta.
\]

由径向与切向方程可得

\[
f=mg\left(\frac{3h}{4l}+\frac32\sin\theta\right),
\qquad
N=\frac34mg\cos\theta.
\]

故首次达到摩擦极限时满足

\[
\frac hl+2\sin\theta_0=\mu\cos\theta_0.
\]

本题 \(h=l\)、\(\mu=2\)
立即恢复式（29）。这条一般式也便于判断某些参数下是否在杆转到竖直前就已经滑动。

\needspace{3\baselineskip}
\textbf{固定支点阶段的时间积分与转向检查}

图6-2左端下降时 \(\theta\) 逆时针增大，球心下降
\(d\sin\theta\)。绕接触点的惯量是 \(I_O=ml^2/9\)，所以

\[
\dot\theta^2=\omega_i^2+\frac{3g}{l}\sin\theta,
\qquad \omega_i^2=\frac{9gh}{2l^2}.
\]

从水平位置走到角度 \(\theta\) 的时间可直接积分：

\[
t(\theta)=\int_0^\theta
\frac{d\psi}{\sqrt{\omega_i^2+(3g/l)\sin\psi}}.
\]

在首个向下阶段分母为正，\(\theta\) 严格增加。支点法向力
\(N=3mg\cos\theta/4\) 随角度减小，所需切向摩擦为
\(f=mg(3h/(4l)+(3/2)\sin\theta)\)，故

\[
\frac{|f|}{N}=
\frac{h/l+2\sin\theta}{\cos\theta}.
\]

\noindent\begin{minipage}{\linewidth}

题给 $h=l$ 且按参考系数 $\mu=2$ 时，临界式是 $1+2\sin\theta=2\cos\theta$；在 $0\le\theta<\pi/2$ 左端增加、右端减小，交点唯一，求得首次滑动约为 $24.3^\circ$。此时 $N>0$，先发生的是切向滑动，并非失去法向接触。

开始滑动后接触点在杆上的物质位置改变，应使用 $x(t)=OC$ 并保留 $2\dot x\dot\theta$ 项。固定支点阶段的能量式不可继续用于滑动阶段，因为摩擦做功；只有杆和桌面的总能量收支仍可追踪。图中的 $d$ 与滑动后的 $x(t)$ 已分别标明。

\end{minipage}

\subsection{六、关键洞察}\label{ux516dux5173ux952eux6d1eux5bdf-5}

本题有三个不同约束阶段：碰撞前自由落体、碰撞后固定点转动、滑动后一般平面运动。每次约束改变，守恒量和自由度都随之改变；正确分段比计算本身更重要。

\subsection{七、易错点}\label{ux4e03ux6613ux9519ux70b9-5}

\begin{enumerate}
\def\labelenumi{\arabic{enumi}.}
\tightlist
\item
  把碰撞过程写成机械能守恒；完全非弹性碰撞中机械能不守恒。
\item
  对质心取角动量守恒，漏掉凸起冲量的力矩。
\item
  把碰撞点到质心距离误写成 \(l/3\)，正确值是 \(l/6\)。
\item
  用 \(f=\mu N\) 描述整个静摩擦阶段；只有临界时才取等号。
\item
  求法向力时遗漏质心切向加速度。
\item
  滑动后仍把凸起当固定转轴，继续使用 \(I_O\ddot\theta=\tau_O\)。
\item
  碰撞后质心速度写成碰前的 \(v_0\)；约束冲量已使速度骤降到 \(v_0/4\)。
\item
  径向方程中把向心项写成 \(-\omega^2d\) 后又在力的方向上重复加负号。
\item
  解半角方程时保留负根，得到与实际逆时针转动方向不符的角度。
\item
  滑动后固定写 \(f=+\mu N\)，没有根据接触点相对滑动速度决定摩擦方向。
\item
  未检查
  \(N\ge0\)；若法向力变为负值，杆已经与凸起分离，接触方程不再适用。
\item
  把附加问的 \(x\) 当成固定的
  \(l/6\)，实际上滑动后接触点在杆上的位置会变化。
\end{enumerate}

\subsection{八、课堂讲解建议}\label{ux516bux8bfeux5802ux8bb2ux89e3ux5efaux8bae-5}

先画一条时间轴，标出``自由下落---碰撞冲量---固定接触转动---开始滑动''四个节点。要求学生在每个区间上方写出可用工具：机械能、角动量、定轴机械能、一般平面运动。这样可以防止把不同阶段的守恒量混用。

碰撞阶段先写角动量式 \(mv_0d=I_O\omega_1\)，随后用
\(V_C=\omega_1d\)。立刻计算碰前、碰后动能，让学生看到机械能损失而角动量守恒。

固定接触阶段在图 6-2 上建立沿杆与垂杆两个方向。先由能量求
\(\omega(\theta)\)，再由径向方程求 \(f\)、切向方程求 \(N\)。不要一开始写
\(f=2N\)；只有两条函数都得到后才在首次交点取临界等号。

滑动后的附加问重点不是求闭式解，而是正确选择广义变量
\(x,\theta\)。让学生从极坐标加速度

\[
a_\parallel=\ddot x-x\dot\theta^2,
\qquad
a_\perp=x\ddot\theta+2\dot x\dot\theta
\]

逐项解释离心项与科里奥利型项，再加转动方程形成闭合系统。最后说明若滑动方向反转或接触丢失，方程还要再次切换。

\subsection{九、变式训练}\label{ux4e5dux53d8ux5f0fux8badux7ec3-5}

把碰撞点改为距质心任意距离 \(d\)。绕碰撞点转动惯量为

\[
I_O=I_C+md^2,
\]

碰撞后角速度与质心速度分别为

\[
\omega_1=\frac{mv_0d}{I_C+md^2},
\qquad
V_C=\omega_1d.
\]

令 \(d\to0\)，角动量臂趋于零，碰撞几乎不激发转动；令 \(d\)
接近杆端，角动量臂增大，但 \(I_O\)
也增大，二者存在竞争。可继续按正文方法反求水平时刻的 \(f,N\)，得到
\(h_{\max}(d,\mu)\)。

这个一般结果可以明确算出。碰后初始时杆仍水平，质心的径向加速度大小为
\(\omega_1^2d\)，因而沿杆摩擦力为

\[
f=m\omega_1^2d.
\]

重力对碰撞点的力矩使角加速度大小为
\(\alpha=mgd/I_O\)。质心沿竖直方向的切向加速度大小为 \(\alpha d\)，故

\[
mg-N=m\alpha d,
\]

从而

\[
N=mg\left(1-\frac{md^2}{I_O}\right)
=mg\frac{I_C}{I_O}.
\]

再代入

\[
\omega_1=\frac{mv_0d}{I_O},
\qquad
v_0^2=2gh,
\]

可得

\[
\frac fN
=\frac{2m^2hd^3}{I_CI_O}.
\]

对均匀细杆 \(I_C=ml^2/12\)，静摩擦条件 \(f/N\le\mu\) 给出

\[
\boxed{
h_{\max}(d,\mu)
=\frac{\mu l^2(l^2/12+d^2)}{24d^3}
}.
\]

令 \(d=l/6\)，立即化为 \(h_{\max}=\mu l\)，与正文完全一致。当 \(d\to0\)
时，碰后转动和径向摩擦都迅速减小，所以允许的高度上限趋于很大；这与上式的
\(d^{-3}\) 趋势相符。

若法向恢复系数为 \(e>0\)，接触点法向速度满足

\[
v_{O,n}^{+}=-e\,v_{O,n}^{-},
\]

而切向是否粘住仍要由摩擦冲量检验。此时``角动量守恒 + 恢复系数 +
切向冲量状态''共同决定碰后速度，不能再直接令接触点速度为零。

\newpage

\section{第 7 题}\label{ux7b2c-7-ux9898}

顶角为 \(45^\circ\) 的金属导轨 \(MON\)
固定在水平面内，并处于竖直方向、磁感应强度为 \(B\)
的匀强磁场中。一根始终垂直于导轨 \(ON\)
的导体棒与两导轨接触，沿水平方向向右运动。导体棒质量为
\(m\)，导轨与导体棒单位长度的电阻均为 \(r_0\)，接触电阻忽略。以顶点
\(O\) 为坐标原点，棒的位置坐标为 \(x\)。

\begin{enumerate}
\def\labelenumi{\arabic{enumi}.}
\tightlist
\item
  棒以恒定速度 \(v_0\) 运动时，求外力 \(F\)。
\item
  求从 \(0\) 到 \(t\) 时间内产生的焦耳热 \(Q(t)\)。
\item
  在 \(t=t_0\) 时撤去外力，求棒最终位置 \(x_f\)。
\item
  若棒从 \(O\) 点以初速度 \(v_1\) 出发，此后不受外力，求最终位置及
  \(x(t)\)。
\end{enumerate}

\begin{figure}[H]
\centering
\includegraphics[width=0.447\linewidth,height=\textheight,keepaspectratio,alt={图7-1　匀强磁场中的直角三角形变长度导轨回路}]{../figures/question-07.pdf}
\caption*{图7-1　匀强磁场中的直角三角形变长度导轨回路}
\end{figure}

\subsection{一、题型定位}\label{ux4e00ux9898ux578bux5b9aux4f4d-6}

棒的位置同时改变三角回路面积和总电阻。先由三条边求电阻，再由面积变化求感应电流；撤去外力后，用速度对位置求导得到极限位置和时间函数。

\subsection{二、知识点}\label{ux4e8cux77e5ux8bc6ux70b9-6}

\begin{enumerate}
\def\labelenumi{\arabic{enumi}.}
\tightlist
\item
  动生电动势 \(\mathcal E=B\,dS/dt\)。
\item
  均匀导线电阻等于单位长度电阻乘总长度。
\item
  安培力可由 \(F_B=BIL\) 或能量平衡求得。
\item
  撤去外力后可用 \(m\dot x\,d\dot x/dx=F(x,\dot x)\) 消去时间。
\item
  一阶非线性方程 \(\dot x=v_1[1-(x/x_f)^2]\) 的分离变量积分。
\end{enumerate}

\subsection{三、物理图像}\label{ux4e09ux7269ux7406ux56feux50cf-6}

当棒位于坐标 \(x\) 时，闭合回路是一个直角三角形。棒长为
\(x\)，下方水平导轨长为 \(x\)，斜边导轨长为 \(\sqrt2x\)。因此回路周长与
\(x\) 成正比，面积与 \(x^2\) 成正比。

\subsection{四、解题入口}\label{ux56dbux89e3ux9898ux5165ux53e3-6}

先用纯几何关系同时写出面积 \(S(x)\) 和电阻
\(R(x)\)。再由磁通量求电动势。不要把安培力简单写成 \(BIl\)
后停下：必须指出只有导体棒所受安培力具有水平方向的合力，其有效长度就是棒长
\(x\)。

\subsection{五、解析}\label{ux4e94ux89e3ux6790-6}

\begin{figure}[H]
\centering
\includegraphics[width=0.877\linewidth,height=\textheight,keepaspectratio,alt={图7-2　面积、电动势、电阻、电流与安培力的递推关系}]{../figures/solution-07.pdf}
\caption*{图7-2　面积、电动势、电阻、电流与安培力的递推关系}
\end{figure}

\needspace{3\baselineskip}
\textbf{（一）几何量、电阻与感应电流}

由于顶角为 \(45^\circ\)，导体棒垂直于 \(ON\)，故棒长

\[
L_{ab}=x,
\tag{1}
\]

水平边长

\[
L_{Ob}=x,
\tag{2}
\]

斜边长

\[
L_{Oa}=\sqrt{x^2+x^2}=\sqrt2x.
\tag{3}
\]

回路总长度

\[
L_{\rm all}=(2+\sqrt2)x,
\tag{4}
\]

所以总电阻

\[
R(x)=(2+\sqrt2)r_0x.
\tag{5}
\]

三角形面积

\[
S(x)=\frac12x^2.
\tag{6}
\]

取棒速度为 \(\dot x\)，磁通量大小为

\[
\Phi=BS=\frac12Bx^2.
\tag{7}
\]

感应电动势大小

\[
\begin{aligned}
\mathcal E
&=\left|\frac{d\Phi}{dt}\right|\\
&=Bx\dot x.
\end{aligned}
\tag{8}
\]

因此感应电流

\[
\begin{aligned}
I
&=\frac{\mathcal E}{R}\\
&=\frac{Bx\dot x}{(2+\sqrt2)r_0x}\\
&=\boxed{\frac{B\dot x}{(2+\sqrt2)r_0}}.
\end{aligned}
\tag{9}
\]

位置因子 \(x\) 被约去，这是后续计算简化的根源。

\needspace{3\baselineskip}
\textbf{（二）匀速运动所需外力}

导体棒所受安培力大小

\[
F_B=BIL_{ab}=BIx.
\tag{10}
\]

由楞次定律，安培力总与棒的运动方向相反。匀速时外力大小等于安培力：

\[
F=F_B.
\tag{11}
\]

将式（9）中 \(\dot x=v_0\) 代入：

\[
\boxed{
F(x)=\frac{B^2xv_0}{(2+\sqrt2)r_0}
}.
\tag{12}
\]

若写成时间函数，由 \(x=v_0t\)：

\[
\boxed{
F(t)=\frac{B^2v_0^2t}{(2+\sqrt2)r_0}
}.
\tag{13}
\]

\needspace{3\baselineskip}
\textbf{（三）匀速阶段产生的焦耳热}

匀速阶段动能不变，外力做功全部转化为焦耳热：

\[
Q(t)=\int_0^tF(t')v_0\,dt'.
\tag{14}
\]

代入式（13）：

\[
\begin{aligned}
Q(t)
&=\int_0^t
\frac{B^2v_0^3t'}{(2+\sqrt2)r_0}\,dt'\\
&=\frac{B^2v_0^3}{(2+\sqrt2)r_0}
\left[\frac{t'^2}{2}\right]_0^t\\
&=\boxed{
\frac{B^2v_0^3t^2}{2(2+\sqrt2)r_0}
}.
\end{aligned}
\tag{15}
\]

也可用 \(I^2R\) 积分复核。匀速时

\[
I=\frac{Bv_0}{(2+\sqrt2)r_0},
\qquad
R(t')=(2+\sqrt2)r_0v_0t'.
\]

所以

\[
\begin{aligned}
Q(t)
&=\int_0^t I^2R(t')\,dt'\\
&=\int_0^t
\frac{B^2v_0^3t'}{(2+\sqrt2)r_0}\,dt'\\
&=\frac{B^2v_0^3t^2}{2(2+\sqrt2)r_0},
\end{aligned}
\]

与外力做功法完全一致。这个复核也说明热功率随时间线性增大，而不是常量。

\needspace{3\baselineskip}
\textbf{（四）在 \(t_0\) 时撤去外力}

撤力瞬间

\[
x_0=v_0t_0,
\qquad
\dot x_0=v_0.
\tag{16}
\]

撤力后运动方程为

\[
m\frac{d\dot x}{dt}
=-\frac{B^2x\dot x}{(2+\sqrt2)r_0}.
\tag{17}
\]

在 \(\dot x>0\) 的阶段使用链式法则

\[
\frac{d\dot x}{dt}
=\frac{d\dot x}{dx}\frac{dx}{dt}
=\dot x\frac{d\dot x}{dx}.
\tag{18}
\]

代入式（17）并约去 \(\dot x\)：

\[
m\frac{d\dot x}{dx}
=-\frac{B^2x}{(2+\sqrt2)r_0}.
\tag{19}
\]

从 \((x_0,v_0)\) 积分到 \((x,\dot x)\)：

\[
\begin{aligned}
\dot x-v_0
&=-\frac{B^2}{(2+\sqrt2)mr_0}
\int_{x_0}^{x}x'\,dx'\\
&=-\frac{B^2}{2(2+\sqrt2)mr_0}
(x^2-x_0^2).
\end{aligned}
\tag{20}
\]

最终速度趋于零，令 \(x=x_f\)、\(\dot x=0\)：

\[
v_0=\frac{B^2}{2(2+\sqrt2)mr_0}(x_f^2-x_0^2).
\tag{21}
\]

所以

\[
\boxed{
x_f=\sqrt{(v_0t_0)^2+\frac{2(2+\sqrt2)mr_0v_0}{B^2}}
}.
\tag{22}
\]

\needspace{3\baselineskip}
\textbf{（五）从顶点以初速度 \(v_1\) 出发}

此时初始条件为

\[
x(0)=0,
\qquad
\dot x(0)=v_1.
\tag{23}
\]

沿用式（19），从 \(0\) 积分到 \(x\)：

\[
\dot x=v_1-\frac{B^2x^2}{2(2+\sqrt2)mr_0}.
\tag{24}
\]

令 \(\dot x\to0\)，得到极限位置

\[
\boxed{
x_f=\sqrt{\frac{2(2+\sqrt2)mr_0v_1}{B^2}}
}.
\tag{25}
\]

利用式（25）把式（24）改写为

\[
\frac{dx}{dt}
=v_1\left(1-\frac{x^2}{x_f^2}\right).
\tag{26}
\]

分离变量：

\[
\frac{dx}{1-x^2/x_f^2}=v_1\,dt.
\tag{27}
\]

从 \(0\) 积分到 \(x(t)\)：

\[
\frac{x_f}{2}
\ln\frac{x_f+x}{x_f-x}=v_1t.
\tag{28}
\]

等价地

\[
\operatorname{artanh}\frac{x}{x_f}
=\frac{v_1t}{x_f}.
\tag{29}
\]

因此

\[
\boxed{
x(t)=x_f\tanh\left(\frac{v_1t}{x_f}\right)
}.
\tag{30}
\]

当 \(t\to\infty\)
时，\(x\to x_f\)、\(\dot x\to0\)。严格地说，棒以指数方式渐近停止，而不是在有限时间突然静止。

\needspace{3\baselineskip}
\textbf{电流方向与一般撤力时刻的显式解}

图7-1取磁场向纸内。棒向右时，面积增大，感应电流沿 \(O\to b\to a\to O\)
逆时针流动，棒上电流向上， \(I\boldsymbol L\times\boldsymbol B\)
指向左，阻碍运动。

记 \(\alpha=B^2/[(2+\sqrt2)mr_0]\)，
\(x_f^2=x_0^2+2v_0/\alpha\)。撤力后的速度可写成

\[
\dot x=\frac\alpha2(x_f^2-x^2).
\]

从 \((t_0,x_0)\) 积分得

\[
\boxed{x(t)=x_f\tanh\left[
\frac{\alpha x_f}{2}(t-t_0)
+\operatorname{artanh}\frac{x_0}{x_f}\right].}
\]

在 \(t=t_0\)，它恢复 \(x_0\)；求导恢复 \(v_0\)。当 \(x_0=0\) 且初速度为
\(v_1\) 时，\(\alpha x_f/2=v_1/x_f\)，便回到式（30）。

图7-2用归一化的实际双曲正切曲线画出 \(x/x_f\) 和 \(v/v_1\)，
并把机械功率与焦耳热功率相等标在同图。回路恰在 \(x=0\) 时退化， 正文采用
\(x\to0^+\) 的理想连续极限，并忽略电感与接触电阻；
若加入有限接触电阻，初始电流极限和运动曲线都需要重算。

\needspace{3\baselineskip}
\textbf{补充推导与适用条件}

\textbf{安培力方向的两种判断。}
楞次定律说明感应电流产生的磁通总要阻碍面积增加，因此棒受力向左。也可用功率核验：若取阻力大小
\(F_B=BIx\)，则

\[
F_B\dot x
=BIx\dot x
=I\mathcal E
=I^2R>0.
\]

机械能减少率恰等于焦耳热功率。若安培力方向写成与运动同向，能量关系会立刻出现矛盾。

\textbf{撤力后的第一积分。} 定义

\[
\alpha=\frac{B^2}{(2+\sqrt2)mr_0}.
\]

运动方程为 \(\ddot x=-\alpha x\dot x\)。在 \(\dot x>0\) 时约去
\(\dot x\) 后得到

\[
\dot x+\frac\alpha2x^2=\text{常量}.
\]

这不是机械能守恒式，而是利用特殊速度型阻力得到的速度---位置第一积分。常量由撤力瞬间或初速度条件确定。

\textbf{为什么只在无限时间停止。} 第四问的解

\[
x=x_f\tanh(v_1t/x_f)
\]

给出

\[
x_f-x\sim2x_f e^{-2v_1t/x_f}\qquad(t\to\infty).
\]

因此剩余距离指数衰减，任何有限时刻仍有极小的正速度。题目所说``最终静止位置''应理解为渐近极限位置。

\textbf{量纲检查。} 单位长度电阻 \(r_0\) 的量纲为
\(\Omega\,\mathrm m^{-1}\)。组合 \(mr_0v_1/B^2\)
的量纲为长度平方，故式（25）根号后确为长度；焦耳热式中
\(B^2v_0^3t^2/r_0\) 的量纲为能量。

\textbf{撤力后的能量核验。} 记

\[
k_B=\frac{B^2}{(2+\sqrt2)r_0},
\]

则安培阻力大小为 \(F_B=k_Bx\dot x\)。撤力后机械能只有动能，其变化率为

\[
\frac{d}{dt}\left(\frac12m\dot x^2\right)
=m\dot x\ddot x
=-k_Bx\dot x^2.
\]

另一方面，电流与总电阻分别为

\[
I=\frac{B\dot x}{(2+\sqrt2)r_0},
\qquad
R=(2+\sqrt2)r_0x,
\]

所以

\[
I^2R=k_Bx\dot x^2.
\]

于是

\[
-\frac{dK}{dt}=I^2R.
\]

从撤力时刻积分到 \(t\to\infty\)，末动能为零，得到撤力后的总焦耳热恰为

\[
\boxed{Q_{\rm after}=\frac12m\dot x_0^2}.
\]

第（3）中 \(\dot x_0=v_0\)，第（4）中初速为
\(v_1\)，因而它们撤力后的全部初动能最终都转化为焦耳热。这一结论与速度---位置第一积分相互独立，可用于反查阻力系数。

\textbf{用显式解反查初值与加速度。} 对第（4）的解求导：

\[
\dot x(t)
=v_1\operatorname{sech}^2
\left(\frac{v_1t}{x_f}\right).
\]

因此 \(\dot x(0)=v_1\)，且任意有限时刻都有 \(\dot x>0\)。再求导得

\[
\ddot x(t)
=-\frac{2v_1^2}{x_f}
\operatorname{sech}^2
\left(\frac{v_1t}{x_f}\right)
\tanh\left(\frac{v_1t}{x_f}\right)\le0.
\]

它在 \(t=0\) 时为零，是因为初始回路尺寸为零、安培力按 \(x\)
消失；随后阻力先增大，到运动后期又随速度衰减而趋零。组合
\(\tau=x_f/v_1\) 是自然的衰减时间尺度；当 \(t\gg\tau\) 时，速度按
\(4v_1e^{-2t/\tau}\)
衰减。这些性质同时核对了双曲正切解的系数、初值和大时间极限。

\subsection{六、关键洞察}\label{ux516dux5173ux952eux6d1eux5bdf-6}

回路的面积导数给出 \(\mathcal E\propto x\dot x\)，回路总电阻给出
\(R\propto x\)，所以 \(I\propto\dot x\)；再乘棒长 \(x\)，阻力成为
\(F_B\propto x\dot x\)。这个结构使得用速度对位置求导后方程可直接积分。

\subsection{七、易错点}\label{ux4e03ux6613ux9519ux70b9-6}

\begin{enumerate}
\def\labelenumi{\arabic{enumi}.}
\tightlist
\item
  总电阻只计算导体棒，漏掉两段导轨。
\item
  斜导轨长度误写成 \(x\)，正确为 \(\sqrt2x\)。
\item
  面积写成 \(x^2\)，漏掉三角形的 \(1/2\)。
\item
  把匀速外力误认为常量；由于棒长增加，\(F\propto x\)。
\item
  在撤力后直接使用常加速度公式。
\item
  把渐近位置理解成有限时间到达的位置。
\item
  在 \(x=0\) 处看到 \(\mathcal E=R=0\) 就武断地说电流不确定；应取
  \(x\to0^+\) 的比值极限。
\item
  安培力用错长度，把斜导轨的 \(\sqrt2x\) 当成运动棒长度。
\item
  焦耳热积分中把 \(R\) 当常量，漏掉 \(R\propto t\)。
\item
  撤力后约去 \(\dot x\) 时没有说明只在 \(\dot x>0\) 的运动阶段进行。
\item
  第（3）问把撤力前已经走过的 \(x_0=v_0t_0\)
  漏掉，只保留撤力后的附加位移。
\item
  第（4）问在 \(t=0\)
  直接代入有限回路公式而不说明是从极小正位置取连续极限。
\end{enumerate}

\subsection{八、课堂讲解建议}\label{ux516bux8bfeux5802ux8bb2ux89e3ux5efaux8bae-6}

先只做几何：让学生从图 7-1 写出三条边 \(x,x,\sqrt2x\) 和面积
\(x^2/2\)。再建立一条因果链

\[
S(x)\rightarrow\Phi(x)\rightarrow\mathcal E(x,\dot x)
\rightarrow I(\dot x)\rightarrow F_B(x,\dot x).
\]

每得到一步，都先判断关于 \(x\)、\(\dot x\)
的次数，再代系数。这样能清楚看到为何电流不含 \(x\)，而安培力仍正比于
\(x\)。

匀速阶段用两条路线求热量：外力做功与
\(I^2R\)。要求两式在积分前就有相同的瞬时功率。撤力后不要从时间方程硬解，先问``阻力同时含
\(x\) 与 \(\dot x\)，怎样用链式法则约去时间''，引出
\(\ddot x=\dot x\,d\dot x/dx\)。

第四问积分到双曲正切后，用图像说明 \(x(t)\) 单调增加、凹向下并渐近
\(x_f\)。再写大时间指数尾项，解释为什么数学上的停止时刻是无穷远，而实验上会在有限精度下看似停止。

\subsection{九、变式训练}\label{ux4e5dux53d8ux5f0fux8badux7ec3-6}

将顶角改为任意 \(\alpha\)。棒位于横坐标 \(x\)
时，棒长、斜边长、面积分别为

\[
L=x\tan\alpha,
\qquad
L_s=\frac{x}{\cos\alpha},
\qquad
S=\frac12x^2\tan\alpha.
\]

令

\[
C_\alpha=1+\tan\alpha+\sec\alpha,
\]

则

\[
R=r_0xC_\alpha,
\qquad
I=\frac{B\tan\alpha}{r_0C_\alpha}\dot x,
\]

\[
F_B=\frac{B^2\tan^2\alpha}{r_0C_\alpha}x\dot x.
\]

从顶点以 \(v_1\) 出发时

\[
\boxed{
x_f^2=\frac{2mr_0v_1C_\alpha}{B^2\tan^2\alpha}
}.
\]

\(\alpha=45^\circ\) 时
\(C_\alpha=2+\sqrt2\)，恢复正文。若单位长度电阻改为
\(r(s)=r_0(1+\lambda s)\)，总电阻不再严格正比于
\(x\)，电流会重新含位置；可建立新的 \(m\ddot x=-BI(x)L(x)\) 数值方程。

\newpage

\section{第 8 题}\label{ux7b2c-8-ux9898}

存在连续衰变链

\[
{}^{228}\mathrm{Th}\longrightarrow{}^{224}\mathrm{Ra}
\longrightarrow{}^{220}\mathrm{Rn}\longrightarrow\cdots,
\]

前两次衰变的半衰期分别为 \(\tau_1=1.91\,\mathrm{a}\) 和
\(\tau_2=3.66\,\mathrm d\)。某矿物中三种核素已达到放射性平衡。从中提取全部
\({}^{228}\mathrm{Th}\) 与 \({}^{224}\mathrm{Ra}\)，组成质量
\(M=1\,\mathrm g\) 的混合物。求：

\begin{enumerate}
\def\labelenumi{\arabic{enumi}.}
\tightlist
\item
  \(t=0\) 时两种核素的核数 \(N_{10},N_{20}\)；
\item
  \(t>0\) 时 \({}^{224}\mathrm{Ra}\) 的核数 \(N_2(t)\)；
\item
  当 \({}^{228}\mathrm{Th}\) 减为初值一半时，\({}^{224}\mathrm{Ra}\)
  的核数。
\end{enumerate}

\begin{figure}[H]
\centering
\includegraphics[width=0.88\linewidth,height=\textheight,keepaspectratio,alt={图8-1　钍、镭、氡的两级连续放射性衰变链}]{../figures/question-08.pdf}
\caption*{图8-1　钍、镭、氡的两级连续放射性衰变链}
\end{figure}

\subsection{一、题型定位}\label{ux4e00ux9898ux578bux5b9aux4f4d-7}

初始平衡给出钍与镭的活度比，用它和一克的质量条件确定核数。提取后钍继续生成镭，须解带源项的衰变方程，再检查钍减半时镭的滞后修正。

\subsection{二、知识点}\label{ux4e8cux77e5ux8bc6ux70b9-7}

\begin{enumerate}
\def\labelenumi{\arabic{enumi}.}
\tightlist
\item
  衰变常量 \(\lambda=\ln2/T\)。
\item
  放射性平衡时母核与子核的衰变率相等。
\item
  两级衰变方程：\(\dot N_1=-\lambda_1N_1\)，\(\dot N_2=\lambda_1N_1-\lambda_2N_2\)。
\item
  一阶线性方程的积分因子法。
\item
  原子核数与质量的换算。
\end{enumerate}

\subsection{三、物理图像}\label{ux4e09ux7269ux7406ux56feux50cf-7}

提取瞬间，混合物中的钍不断给镭补充核数，镭自身又不断衰变。由于
\(\tau_2\ll \tau_1\)，镭对母核缓慢变化的响应很快：经历几个 \(\tau_2\)
后，初始镭的记忆项基本消失，镭核数近似跟随钍核数缓慢下降。

\subsection{四、解题入口}\label{ux56dbux89e3ux9898ux5165ux53e3-7}

先由平衡态得到核数比，再与质量方程联立确定初值。求 \(N_2(t)\)
时必须同时保留``母核持续生成''与``子核自身衰变''两项；不能把镭简单写成单独的指数衰减。

\subsection{五、解析}\label{ux4e94ux89e3ux6790-7}

\begin{figure}[H]
\centering
\includegraphics[width=0.88\linewidth,height=\textheight,keepaspectratio,alt={图8-2　母核与短寿命子核在准平衡条件下的时间演化}]{../figures/solution-08.pdf}
\caption*{图8-2　母核与短寿命子核在准平衡条件下的时间演化}
\end{figure}

\needspace{3\baselineskip}
\textbf{（一）初始核数}

定义衰变常量

\[
\lambda_1=\frac{\ln2}{\tau_1},
\qquad
\lambda_2=\frac{\ln2}{\tau_2}.
\tag{1}
\]

矿物中已达到平衡态，意味着每单位时间由钍生成的镭核数等于镭的衰变核数：

\[
\lambda_1N_{10}=\lambda_2N_{20}.
\tag{2}
\]

因此

\[
\frac{N_{10}}{N_{20}}
=\frac{\lambda_2}{\lambda_1}
=\frac{\tau_1}{\tau_2}.
\tag{3}
\]

统一时间单位。取一年为 \(365.25\) 天：

\[
\tau_1=1.91\times365.25\,\mathrm d
=697.6275\,\mathrm d.
\tag{4}
\]

故

\[
\rho\equiv\frac{\tau_1}{\tau_2}
=\frac{697.6275}{3.66}
\approx190.61.
\tag{5}
\]

混合物总质量为
\(1\,\mathrm g\)。忽略电子质量和质量亏损对本题有效数字的影响，有

\[
\frac{228N_{10}+224N_{20}}{N_A}\,\mathrm g
=1\,\mathrm g,
\tag{6}
\]

其中 \(N_A=6.02214076\times10^{23}\,\mathrm{mol^{-1}}\)。由
\(N_{10}=\rho N_{20}\)：

\[
(228\rho+224)N_{20}=N_A.
\tag{7}
\]

于是

\[
\begin{aligned}
N_{20}
&=\frac{N_A}{228\rho+224}\\
&\approx\boxed{1.38\times10^{19}},
\end{aligned}
\tag{8}
\]

以及

\[
\begin{aligned}
N_{10}
&=\rho N_{20}\\
&\approx\boxed{2.63\times10^{21}}.
\end{aligned}
\tag{9}
\]

核数比约为
\(191\)，所以尽管单个钍核略重，混合物质量几乎全部来自钍；数量级结果与此直觉一致。

\needspace{3\baselineskip}
\textbf{（二）\(t>0\) 时的镭核数}

母核满足

\[
\frac{dN_1}{dt}=-\lambda_1N_1,
\tag{10}
\]

所以

\[
N_1(t)=N_{10}e^{-\lambda_1t}.
\tag{11}
\]

镭核一方面由钍衰变产生，产生率为 \(\lambda_1N_1\)；另一方面自身以速率
\(\lambda_2N_2\) 衰变。因此

\[
\frac{dN_2}{dt}
=\lambda_1N_{10}e^{-\lambda_1t}-\lambda_2N_2.
\tag{12}
\]

整理为标准线性形式：

\[
\frac{dN_2}{dt}+\lambda_2N_2
=\lambda_1N_{10}e^{-\lambda_1t}.
\tag{13}
\]

乘积分因子 \(e^{\lambda_2t}\)：

\[
e^{\lambda_2t}\frac{dN_2}{dt}
+\lambda_2e^{\lambda_2t}N_2
=\lambda_1N_{10}e^{(\lambda_2-\lambda_1)t}.
\tag{14}
\]

左边正好是乘积导数：

\[
\frac{d}{dt}\left(N_2e^{\lambda_2t}\right)
=\lambda_1N_{10}e^{(\lambda_2-\lambda_1)t}.
\tag{15}
\]

从 \(0\) 到 \(t\) 积分：

\[
N_2(t)e^{\lambda_2t}-N_{20}
=\lambda_1N_{10}
\frac{e^{(\lambda_2-\lambda_1)t}-1}
{\lambda_2-\lambda_1}.
\tag{16}
\]

两边乘 \(e^{-\lambda_2t}\)：

\[
\boxed{
N_2(t)=
\frac{\lambda_1N_{10}}{\lambda_2-\lambda_1}
\left(e^{-\lambda_1t}-e^{-\lambda_2t}\right)
+N_{20}e^{-\lambda_2t}
}.
\tag{17}
\]

把 \(t=0\) 代入，第一项为零、第二项为 \(N_{20}\)，初值检验通过。

利用初始平衡关系
\(N_{20}=\lambda_1N_{10}/\lambda_2\)，式（17）也可整理成

\[
N_2(t)=
\frac{\lambda_1N_{10}}{\lambda_2-\lambda_1}
e^{-\lambda_1t}
-\frac{\lambda_1^2N_{10}}
{\lambda_2(\lambda_2-\lambda_1)}e^{-\lambda_2t}.
\tag{18}
\]

第二项系数还多含因子 \(\lambda_1/\lambda_2\approx1/190.61\)，且随
\(e^{-\lambda_2t}\) 很快衰减。

对解作两个动力学检查。首先在 \(t=0\) 求导：

\[
\left.\frac{dN_2}{dt}\right|_{t=0}
=\lambda_1N_{10}-\lambda_2N_{20}=0,
\]

这正是``提取瞬间仍处于平衡配比''的含义。其次，当 \(t\gg \tau_2\)
时，快变项 \(e^{-\lambda_2t}\) 消失，故

\[
N_2(t)\approx
\frac{\lambda_1}{\lambda_2-\lambda_1}N_1(t).
\]

于是长期准平衡核数比为

\[
\frac{N_2}{N_1}
\approx\frac{\lambda_1}{\lambda_2-\lambda_1}
=\frac{\tau_2}{\tau_1-\tau_2},
\]

它比静态比值 \(\lambda_1/\lambda_2\)
略大。这是因为母核数正在缓慢下降，子核对下降存在一个约为 \(\tau_2\)
的响应滞后。

\needspace{3\baselineskip}
\textbf{（三）钍减半时的镭核数}

钍减为初值一半的时刻正是

\[
t=\tau_1,
\qquad
e^{-\lambda_1\tau_1}=\frac12.
\tag{19}
\]

另一方面

\[
e^{-\lambda_2\tau_1}
=2^{-\tau_1/\tau_2}
=2^{-190.61},
\tag{20}
\]

它约为 \(10^{-57}\)，在本题精度下可以忽略。于是由式（17）：

\[
\begin{aligned}
N_2(\tau_1)
&\approx
\frac{\lambda_1N_{10}}{\lambda_2-\lambda_1}
\cdot\frac12\\
&=\frac{N_{20}}{2(1-\lambda_1/\lambda_2)}.
\end{aligned}
\tag{21}
\]

由于 \(\lambda_1/\lambda_2=\tau_2/\tau_1\approx0.005246\)，

\[
\boxed{
N_2(\tau_1)\approx6.93\times10^{18}
}.
\tag{22}
\]

若只取最低阶近似，则 \(N_2(\tau_1)\approx N_{20}/2\)；式（22）保留了
\(\lambda_1\) 相对 \(\lambda_2\) 的一阶修正。

\needspace{3\baselineskip}
\textbf{归一化曲线与平衡含义}

令 \(s=t/\tau_1\)、\(\rho=\tau_1/\tau_2\)。利用初始活度相等，

\[
\frac{N_1(t)}{N_{10}}=2^{-s},\qquad
\frac{N_2(t)}{N_{20}}
=\frac{\rho}{\rho-1}2^{-s}
-\frac1{\rho-1}2^{-\rho s}.
\]

图8-2由这两个函数生成，而不是画两条相同的指数。在原点， 钍的斜率为
\(-\ln2\)，镭的斜率为零；很快镭曲线跟随慢指数，
但归一化核数比略大于钍。图中增加前若干天的放大面板，
以便看清初始水平切线和两曲线的细小差异。

若只有这条链而母核无外部补充，两种核数不可能永久不随时间变化。
原题所说矿物中核数已平衡，按参考解理解为提取时
\(\lambda_1N_{10}=\lambda_2N_{20}\)；持续补充母核的上游来源未写出。
提取后停止该补充，才使用式（10）。不要把``初始活度相等''错误地
当成提取后始终精确满足同一静态比例。

直接代入数值给出更细的回查： \(N_{10}\approx2.62775\times10^{21}\)、
\(N_{20}\approx1.37861\times10^{19}\)（末位按计算值取舍）， 质量
\(228N_{10}/N_A+224N_{20}/N_A\) 应为1克。
镭从零开始长入的变式还可由卷积积分求得同一峰值时刻，
它与本题初始活度平衡时立即缓慢下降的曲线不同。

\needspace{3\baselineskip}
\textbf{补充推导与适用条件}

\textbf{初始平衡为何给出活度相等。} 在矿物中若 \({}^{224}\mathrm{Ra}\)
的平均核数长期不变，则

\[
\frac{dN_2}{dt}=\lambda_1N_1-\lambda_2N_2\approx0.
\]

所以相等的是单位时间衰变数
\(\lambda N\)，不是核数本身。短寿命核素要用较少的核数提供与长寿命母核相同的活度，这正解释了
\(N_{20}\ll N_{10}\)。

\textbf{精确解的两项含义。}
式（17）第一项是提取后由钍不断生成的新镭，第二项是初始已有镭的记忆项。后者按
\(e^{-\lambda_2t}\) 在几周内迅速消失；前者则主要按 \(e^{-\lambda_1t}\)
在数年尺度变化。把两项分开看，比把公式整体背下来更容易判断极限。

\textbf{第三问不是严格减半。} 在 \(t=\tau_1\)
时，快变项虽然小到可忽略，但慢变项系数是

\[
\frac{1}{1-\lambda_1/\lambda_2}>1.
\]

因此 \(N_2(\tau_1)\) 略大于
\(N_{20}/2\)。子核对母核下降存在迟滞，方向上必须是``略大''，若计算出略小则通常是分母符号写反。

\textbf{数量级复核。} 一克、摩尔质量约 \(228\,\mathrm{g\,mol^{-1}}\)
的物质含核数约

\[
\frac{N_A}{228}\approx2.64\times10^{21}.
\]

本题
\(N_{10}\approx2.63\times10^{21}\)，说明混合物质量几乎全由钍贡献；再除以约
\(191\) 得
\(N_{20}\sim10^{19}\)。不借助微分方程也能先预判两个答案的数量级。

\textbf{用``生成时刻''重建子核解。} 在时刻 \(\tau\) 到 \(\tau+d\tau\)
内，母核产生的子核数为

\[
dN_{2,\rm prod}=\lambda_1N_{10}e^{-\lambda_1\tau}d\tau.
\]

这一小批子核要从 \(\tau\) 存活到 \(t\)，存活概率为
\(e^{-\lambda_2(t-\tau)}\)。再加上初始已有子核的存活部分：

\[
N_2(t)=N_{20}e^{-\lambda_2t}
+\int_0^t\lambda_1N_{10}e^{-\lambda_1\tau}
e^{-\lambda_2(t-\tau)}d\tau.
\]

把与 \(\tau\) 无关的因子提出，逐步积分：

\[
\begin{aligned}
N_2(t)
&=N_{20}e^{-\lambda_2t}
+\lambda_1N_{10}e^{-\lambda_2t}
\int_0^t e^{(\lambda_2-\lambda_1)\tau}d\tau\\
&=N_{20}e^{-\lambda_2t}
+\frac{\lambda_1N_{10}}{\lambda_2-\lambda_1}
e^{-\lambda_2t}
\left[e^{(\lambda_2-\lambda_1)t}-1\right]\\
&=N_{20}e^{-\lambda_2t}
+\frac{\lambda_1N_{10}}{\lambda_2-\lambda_1}
\left(e^{-\lambda_1t}-e^{-\lambda_2t}\right).
\end{aligned}
\]

这与积分因子法完全一致，但每个指数因子都有``何时生成、存活多久''的直接意义。

\textbf{提取初期的曲线形状。} 因初始活度平衡，

\[
\dot N_2(0)=\lambda_1N_{10}-\lambda_2N_{20}=0.
\]

再对速率方程求导，并使用 \(\dot N_1(0)=-\lambda_1N_{10}\)：

\[
\ddot N_2(0)
=\lambda_1\dot N_1(0)-\lambda_2\dot N_2(0)
=-\lambda_1^2N_{10}<0.
\]

因而 \(N_2(t)\) 在 \(t=0\)
处切线水平，但随后向下弯曲，并非保持一段严格平台。这个局部展开能检查图
8-2 在原点附近的形状是否合理。

\textbf{快慢时间尺度的层次。} 系统同时含有 \(\tau_2=3.66\,\mathrm d\) 和
\(\tau_1\approx698\,\mathrm d\)
两个尺度。在提取后的前几周，\(e^{-\lambda_2t}\)
的快变记忆项完成调整，而母核只改变很小的比例；在 \(t\gg \tau_2\) 但
\(t\) 仍与 \(\tau_1\)
同阶时，只剩随母核衰减的慢模态。因此可写出准平衡近似

\[
N_2(t)
\approx\frac{\lambda_1}{\lambda_2-\lambda_1}N_1(t)
=\frac{1}{1-\tau_2/\tau_1}
\frac{\lambda_1}{\lambda_2}N_1(t).
\]

与静态平衡比值相比，修正因子为
\(1/(1-\tau_2/\tau_1)\approx1.00527\)，只高约
\(0.527\%\)。这说明第（3）可以用``约为初值一半''快速估算，但在标准解中应保留这个可定量的滞后修正。

\subsection{六、关键洞察}\label{ux516dux5173ux952eux6d1eux5bdf-7}

所谓平衡并不表示每一种核素``停止变化''，而表示在当时生成率与衰变率相等。提取后，系统仍遵循连续衰变方程；由于子核半衰期很短，它会迅速进入跟随母核的准平衡状态。

\subsection{七、易错点}\label{ux4e03ux6613ux9519ux70b9-7}

\begin{enumerate}
\def\labelenumi{\arabic{enumi}.}
\tightlist
\item
  把平衡条件写成 \(N_{10}=N_{20}\)，而不是活度相等。
\item
  质量方程中把两种核素都按同一摩尔质量计算。
\item
  求 \(N_2(t)\) 时遗漏源项 \(\lambda_1N_1\)。
\item
  在指数函数中混用年和天。
\item
  看到钍减半便直接断言镭严格减半；这只是 \(\lambda_2\gg\lambda_1\)
  下的近似。
\item
  把半衰期直接当作衰变常量，漏掉 \(\lambda=\ln2/T\)。
\item
  质量方程忘记除以阿伏伽德罗常数，或把核数直接与克数相加。
\item
  解一阶方程时积分因子符号写反，导致解中出现随时间增长的
  \(e^{+\lambda_2t}\)。
\item
  写出通解后不代 \(t=0\)，没有发现积分常数或两指数系数错误。
\item
  忽略 \(e^{-\lambda_2\tau_1}\) 时只说``很小''，没有用
  \(2^{-\tau_1/\tau_2}\) 说明其约为 \(10^{-57}\)。
\item
  把每个母核衰变都必然产生被保留的子核当成无条件事实；本题连续衰变链默认分支比为
  \(1\)。
\end{enumerate}

\subsection{八、课堂讲解建议}\label{ux516bux8bfeux5802ux8bb2ux89e3ux5efaux8bae-7}

先用``上游水箱缓慢供水、下游水箱快速漏水''解释活度平衡。要求学生用一句话区分``核数相等''和``流率相等''，再写
\(\lambda_1N_{10}=\lambda_2N_{20}\)。

初值计算分两列：左列是核数比，右列是质量方程。统一把时间换成天后再求比值，避免在指数中混用单位。算完先用
\(N_A/228\) 估算总数量级，确认钍核数接近 \(2.6\times10^{21}\)。

微分方程阶段把

\[
\dot N_2=\underbrace{\lambda_1N_1}_{\text{生成}}
-\underbrace{\lambda_2N_2}_{\text{损失}}
\]

写成带颜色的两项。积分因子求解完成后，立即进行三项检查：\(t=0\)
恢复初值；初始导数为零；\(t\gg \tau_2\)
后只剩慢指数。第三问再利用半衰期比约 \(191\)
舍弃快项，并解释答案为何略大于严格的一半。

最后把图 8-2 中两条曲线在早期几乎重合、在 \(\tau_1\)
处同时约为一半的形状与解析式对应起来。

\subsection{九、变式训练}\label{ux4e5dux53d8ux5f0fux8badux7ec3-7}

若提取时只保留钍而完全除去镭，即 \(N_2(0)=0\)，则

\[
N_2(t)=\frac{\lambda_1N_{10}}{\lambda_2-\lambda_1}
\left(e^{-\lambda_1t}-e^{-\lambda_2t}\right).
\]

令导数为零：

\[
\lambda_2e^{-\lambda_2t_*}
=\lambda_1e^{-\lambda_1t_*},
\]

故峰值时刻为

\[
\boxed{
t_*=\frac{\ln(\lambda_2/\lambda_1)}
{\lambda_2-\lambda_1}
}.
\]

该结果体现子核先积累后衰减的``长入曲线''。若再加入
\({}^{220}\mathrm{Rn}\)，则

\[
\dot N_3=\lambda_2N_2-\lambda_3N_3,
\]

把已求出的两个指数代入后会再产生第三个时间尺度。还可加入小于 \(1\)
的分支比 \(b\)，把生成项改为
\(b\lambda_1N_1\)，讨论平衡活度如何随分支比降低。
\end{document}
